% Generated mechanically from frozen input inputs/p08/ja/relative-cycles.tex.
% Do not edit this generated file; edit the bound locale source instead.

% OK, start here.
%

\setcounter{chapter}{61}
\chapter{相対サイクル}



\phantomsection
\label{relative-cycles-section-phantom}




\section{序論}
\label{relative-cycles-section-introduction}

\noindent
基礎的な文献は \cite{SV} である。

\medskip\noindent
本章では、\cite{SV} で普遍整相対サイクルと呼ばれるものだけを定義する。
この選択により、理論の展開は原論文よりいくらか簡単になるが、もちろん失われるものもある。

\medskip\noindent
Noether スキーム間の有限型射 $X \to S$ を固定する。$X/S$ のファイバー上の
$r$-サイクルの族 $\alpha$ とは、単に $\alpha_s \in Z_r(X_s)$ なる集まり
$\alpha = (\alpha_s)_{s \in S}$ のことである。任意の射 $g : S' \to S$ に沿う
$\alpha$ の基底変換 $g^*\alpha$ の定め方は直ちに明らかである。
$\alpha$ が特殊化と両立するとき、すなわち $S'$ が離散付値環のスペクトルである
任意の射 $g : S' \to S$ に対し、$g^*\alpha$ の一般ファイバーが
$g^*\alpha$ の閉ファイバーへ特殊化することを要求するとき、$\alpha$ を
{\it $X/S$ 上の相対 $r$-サイクル}という。Section
\ref{relative-cycles-section-families-specialization} を参照せよ。






\section{規約と記法}
\label{relative-cycles-section-conventions}

\noindent
固定した Noether 基底上局所有限型なスキームのサイクルに関する規約と記法については、
Chow Homology and Chern Classes の章、特に Chow Homology, Section
\ref{chow-section-setup} 以降を参照せよ。

\medskip\noindent
特に、$X$ が体 $k$ 上局所有限型なら、$Z_r(X)$ は次元 $r$ のサイクル群を表す。
Chow Homology, Example \ref{chow-example-field} および Section
\ref{chow-section-cycles} を参照せよ。$\dim(Z) = r$ を満たす整閉部分スキーム
$Z \subset X$ に対して $[Z] \in Z_r(X)$ であり、$X$ が準コンパクトなら
$Z_r(X)$ はこれらの類を基底とする自由アーベル群である。




\section{体に対する相対サイクル}
\label{relative-cycles-section-relative-fields}

\noindent
体 $k$ と、$k$ 上局所代数的なスキーム $X$ を取る。$r \geq 0$ を整数とする。
この状況では、$X$ 上の $r$-サイクルの群 $Z_r(X)$ がある。Section
\ref{relative-cycles-section-conventions} を参照せよ。

\medskip\noindent
{\bf 基底変換。} 任意の体拡大 $k'/k$ に対して基底変換写像
$Z_r(X) \to Z_r(X_{k'})$ がある。Chow Homology, Section
\ref{chow-section-change-base} を参照せよ。具体的には、次元 $r$ の整閉部分スキーム
$Z \subset X$ に対し、$[Z] \in Z_r(X)$ を閉部分スキーム
$Z_{k'} \subset X_{k'}$ に付随する $r$-サイクル
$[Z_{k'}]_r \in Z_r(X_{k'})$ へ送る（もちろん一般には $Z_{k'}$ は
既約でも被約でもない）。基底変換写像 $Z_r(X) \to Z_r(X_{k'})$ は常に単射である。

\begin{lemma}
\label{relative-cycles-lemma-multiplicities-field-extension}
$K/k$ を体拡大とし、$Z$ を $k$ 上の整局所代数的スキームとする。
既約成分 $Z' \subset Z_K$ の重複度 $m_{Z', Z_K}$ は $1$、または
$k$ の標数の冪である。
\end{lemma}

\begin{proof}
$k$ の標数が零なら $k$ は完全であり、Varieties, Lemma
\ref{varieties-lemma-geometrically-reduced} により $X_K$ は被約なので、
重複度は常に $1$ である。$k$ の標数が $p > 0$ であると仮定する。
$L$ を $Z$ の関数体とする。$Z$ は $k$ 上局所代数的なので、体拡大
$L/k$ は有限生成である。環 $K \otimes_k L$ は Noether である
（Algebra, Lemma \ref{algebra-lemma-Noetherian-field-extension}）。
代数の言葉に直すと、任意の極小素イデアル $\mathfrak q$ に対し、Artin 局所環
$(K \otimes_k L)_\mathfrak q$ の長さが $p$ の冪であることを示せばよい。

\medskip\noindent
$L'/L$ を次数 $p^n$ の有限純非分離拡大とする。このとき
$K \otimes_k L \subset K \otimes_k L'$ は次数 $p^n$ の有限自由環準同型であり、
スペクトルの同相写像と純非分離な剰余体拡大を誘導する。したがって、上の任意の
極小素イデアル $\mathfrak q$ に対し、その上にある極小素イデアル
$\mathfrak q' \subset K \otimes_k L'$ はただ一つであり、
$$
p^n \text{length}((K \otimes_k L)_\mathfrak q) =
[\kappa(\mathfrak q') : \kappa(\mathfrak q)]
\text{length}((K \otimes_k L')_{\mathfrak q'})
$$
が成り立つ。これは $M = (K \otimes_k L')_{\mathfrak q'} \cong
(K \otimes_k L)_{\mathfrak q}^{\oplus p^n}$ に Algebra, Lemma
\ref{algebra-lemma-pushdown-module} を適用したものである。
$[\kappa(\mathfrak q') : \kappa(\mathfrak q)]$ は $p$ の冪なので、
$L'$ と $\mathfrak q'$ について主張を証明すれば十分である。

\medskip\noindent
前段落と Algebra, Lemma \ref{algebra-lemma-make-separable} により、
$L/k'$ が分離的で $k'/k$ が有限純非分離となる部分体の塔 $L/k'/k$ があると
仮定してよい。このとき $K \otimes_k k'$ は Artin 局所環である。前段落の議論を
$L = k$ と $L' = k'$ に適用すると、$\text{length}(K \otimes_k k')$ は
$p$ の冪である。$L/k'$ は滑らかな $k'$-代数の局所化なので
（Algebra, Lemma \ref{algebra-lemma-localization-smooth-separable}）、
$S = (K \otimes_k L)_\mathfrak q$ は滑らかな
$R = K \otimes_k k'$-代数を極小素イデアルで局所化したものである。
したがって $R \to S$ は Artin 局所環の平坦な局所準同型であり、
$\mathfrak m_R S = \mathfrak m_S$ である。Algebra, Lemma
\ref{algebra-lemma-pullback-module} により
$\text{length}(K \otimes_k k') = \text{length}(R) =
\text{length}(S) = \text{length}((K \otimes_k L)_\mathfrak q)$
を得て、証明が完了する。
\end{proof}

\begin{lemma}
\label{relative-cycles-lemma-how-different}
$k$ を標数 $p > 0$ の体とし、その完全閉包を $k^{perf}$ とする。
$X$ を $k$ 上の代数的スキーム、$r \geq 0$ を整数とする。単射
$Z_r(X) \to Z_r(X_{k^{perf}})$ の余核は $p$-冪捩れ加群である
（More on Algebra, Definition
\ref{more-algebra-definition-f-power-torsion}）。
\end{lemma}

\begin{proof}
$X$ は準コンパクトなので、アーベル群 $Z_r(X)$ は次元 $r$ の整閉部分スキームを
基底とする自由群である。$Z_r(X_{k^{perf}})$ についても同様である。
$X_{k^{perf}} \to X$ は同相写像なので、$Z_r(X) \to Z_r(X_{k^{perf}})$ は
余核が捩れ群となる単射である。Lemma
\ref{relative-cycles-lemma-multiplicities-field-extension} により、余核の各元は $p$-冪捩れである。
\end{proof}





\section{サイクルの特殊化}
\label{relative-cycles-section-specialization}

\noindent
$R$ を分数体 $K$、剰余体 $\kappa$ をもつ離散付値環とする。
$X$ を $R$ 上局所有限型なスキームとし、$r \geq 0$ とする。特殊化写像
$$
sp_{X/R} : Z_r(X_K) \longrightarrow Z_r(X_\kappa)
$$
を次のように定義する。次元 $r$ の整閉部分スキーム $Z \subset X_K$ に対し、
$Z \to X$ のスキーム論的像を $\overline{Z}$ で表す。そして $sp_{X/R}$ を
$$
sp_{X/R}([Z]) = [\overline{Z}_\kappa]_r
$$
を満たす一意な $\mathbf{Z}$-線形写像とする。これが良定義である理由を簡単に述べる。
まず、射 $X_K \to X$ は準コンパクトであり、したがって射 $Z \to X$ も
準コンパクトである。ゆえに Morphisms, Lemma
\ref{morphisms-lemma-flat-base-change-scheme-theoretic-image} により、
$Z \to X$ のスキーム論的像を取る操作は平坦基底変換と可換する。特に、再び
$X_K$ へ基底変換すると $Z = \overline{Z}_K$ を得る。$Z$ は整なので
$\overline{Z}$ も整であり、実際、$Z$ の一般点の像を一般点とする一意な整閉部分スキームに
等しい。Varieties, Lemma
\ref{varieties-lemma-dominate-valuation-ring-dimension-fibres} により、
$Z_\kappa$ は次元 $r$ の純次元スキームである。

\begin{lemma}
\label{relative-cycles-lemma-specialization-module}
$R$ を分数体 $K$、剰余体 $\kappa$ をもつ離散付値環とする。$X$ を $R$ 上
局所有限型なスキームとし、$r \geq 0$ とする。$\mathcal{F}$ を $R$ 上平坦な
連接 $\mathcal{O}_X$-加群とする。$\dim(\text{Supp}(\mathcal{F}_K)) \leq r$
と仮定する。このとき $\dim(\text{Supp}(\mathcal{F}_\kappa)) \leq r$ であり、
$$
sp_{X/R}([\mathcal{F}_K]_r) = [\mathcal{F}_\kappa]_r
$$
\end{lemma}

\begin{proof}
次元についての主張は More on Morphisms, Lemma
\ref{more-morphisms-lemma-relative-dimension-support-flat} から従う。
$x$ を、次元 $r$ の整閉部分スキーム $Z \subset X_\kappa$ の一般点とする。
証明を完了するため、この等式の左辺 (L) と右辺 (R) における $[Z]$ の係数が
等しいことを示す。

\medskip\noindent
$A = \mathcal{O}_{X, x}$ および $M = \mathcal{F}_x$ とおく。
$M$ は $R$ 上平坦な有限 $A$-加群である。$\pi \in R$ を
$A/\pi A = \mathcal{O}_{X_\kappa, x}$ となる一様化元とする。
Chow Homology, Lemma \ref{chow-lemma-additivity-divisors-restricted} により
$$
\sum\nolimits_i \text{length}_A(A/(\pi, \mathfrak q_i))
\text{length}_{A_{\mathfrak q_i}}(M_{\mathfrak q_i}) =
\text{length}_A(M/\pi M)
$$
を得る。ここで和は $M$ の台に含まれる極小素イデアル $\mathfrak q_i$ にわたる。
$\pi$ は $M$ 上の非零因子なので $\pi \not \in \mathfrak q_i$ であり、したがって
これらの素イデアルは、選んだ $x \in X_\kappa$ へ特殊化する
$\mathcal{F}_K$ の台の一般点 $y_i \in X_K$ に対応する。ゆえに左辺は (L) における
$[Z]$ の係数である。一方、$\text{length}_A(M/\pi M)$ は (R) における
$[Z]$ の係数である。これで証明が完了する。
\end{proof}

\begin{lemma}
\label{relative-cycles-lemma-specialization-closed}
$R$ を分数体 $K$、剰余体 $\kappa$ をもつ離散付値環とする。$X$ を $R$ 上
局所有限型なスキームとし、$r \geq 0$ とする。$W \subset X$ を $R$ 上平坦な
閉部分スキームとする。$\dim(W_K) \leq r$ と仮定する。このとき
$\dim(W_\kappa) \leq r$ であり、
$$
sp_{X/R}([W_K]_r) = [W_\kappa]_r
$$
\end{lemma}

\begin{proof}
$\mathcal{F} = \mathcal{O}_W$ とおけば、これは
Lemma \ref{relative-cycles-lemma-specialization-module} の特別な場合である。Chow Homology,
Lemma \ref{chow-lemma-cycle-closed-coherent} も参照せよ。
\end{proof}

\begin{lemma}
\label{relative-cycles-lemma-specialization-extension}
$R'/R$ を、分数体の拡大 $K'/K$ と剰余体の拡大 $\kappa'/\kappa$ を誘導する
離散付値環の拡大とする (More on Algebra, Definition
\ref{more-algebra-definition-extension-discrete-valuation-rings})。
$X$ を $R$ 上局所有限型とし、$X' = X_{R'}$ と書く。このとき図式
$$
\xymatrix{
Z_r(X'_{K'}) \ar[rr]_{sp_{X'/R'}} & & Z_r(X'_{\kappa'}) \\
Z_r(X_K) \ar[rr]^{sp_{X/R}} \ar[u] & & Z_r(X_\kappa) \ar[u]
}
$$
は可換である。ここで $r \geq 0$ であり、縦の射は基底変換写像である。
\end{lemma}

\begin{proof}
$X'_{K'} = X_{K'} = X_K \times_{\Spec(K)} \Spec(K')$ であり、閉ファイバーに
ついても同様なので、縦の射は実際に定義される
(Section \ref{relative-cycles-section-relative-fields} を参照せよ)。いま $Z \subset X_K$ を、
スキーム論的像が $\overline{Z} \subset X$ である整閉部分スキームとする。このとき
$Z_{K'} \subset X_{K'}$ は、スキーム論的像が
$\overline{Z}_{R'} \subset X_{R'}$ である閉部分スキームである。定義により、
$[Z]$ の基底変換は $[Z_{K'}]_r = [\overline{Z}_{K'}]_r$ である。また
$$
sp_{X/R}([Z]) = [\overline{Z}_\kappa]_r
\quad\text{かつ}\quad
sp_{X'/R'}([\overline{Z}_{K'}]_r) = [(\overline{Z}_{R'})_{\kappa'}]_r
$$
が Lemma \ref{relative-cycles-lemma-specialization-module} により成り立つ。
$(\overline{Z}_{R'})_{\kappa'} = (\overline{Z}_\kappa)_{\kappa'}$ なので結論を得る。
\end{proof}

\begin{lemma}
\label{relative-cycles-lemma-specialization-flat-pullback}
$R$ を分数体 $K$、剰余体 $\kappa$ をもつ離散付値環とする。$X$ を $R$ 上
局所有限型なスキームとする。$f : X' \to X$ を、局所有限型かつ平坦で相対次元
$e$ の射とする。このとき図式
$$
\xymatrix{
Z_{r + e}(X'_K) \ar[rr]_{sp_{X'/R}} & & Z_{r + e}(X'_\kappa) \\
Z_r(X_K) \ar[rr]^{sp_{X/R}} \ar[u] & & Z_r(X_\kappa) \ar[u]
}
$$
は可換である。ここで $r \geq 0$ であり、縦の射は平坦引き戻しによる。
\end{lemma}

\begin{proof}
$Z \subset X$ を $R$ を支配する整閉部分スキームとする。$sp_{X/R}$ の構成により
$sp_{X/R}([Z_K]) = [Z_\kappa]_r$ であり、この性質が特殊化写像を特徴づける。
$Z' = f^{-1}(Z) = X' \times_X Z$ とおく。$R$ は付値環なので $Z$ は $R$ 上
平坦である。したがって $Z'$ も $R$ 上平坦であり、
$sp_{X'/R}([Z'_K]_{r + e}) = [Z'_\kappa]_{r + e}$
が Lemma \ref{relative-cycles-lemma-specialization-closed} により成り立つ。さらに Chow Homology,
Lemma \ref{chow-lemma-pullback-coherent} により
$f_K^*[Z_K] = [Z'_K]_{r + e}$ および
$f_\kappa^*[Z_\kappa]_r = [Z'_\kappa]_{r + e}$ なので、主張を得る。
\end{proof}

\begin{lemma}
\label{relative-cycles-lemma-specialization-proper-pushforward}
$R$ を分数体 $K$、剰余体 $\kappa$ をもつ離散付値環とする。
$f : X \to Y$ を、$R$ 上局所有限型なスキームの固有射とする。このとき図式
$$
\xymatrix{
Z_r(X_K) \ar[rr]_{sp_{X/R}} \ar[d] & & Z_r(X_\kappa) \ar[d] \\
Z_r(Y_K) \ar[rr]^{sp_{Y/R}} & & Z_r(Y_\kappa)
}
$$
は可換である。ここで $r \geq 0$ であり、縦の射は固有前送りによる。
\end{lemma}

\begin{proof}
$Z \subset X$ を $R$ を支配する整閉部分スキームとする。$sp_{X/R}$ の構成により
$sp_{X/R}([Z_K]) = [Z_\kappa]_r$ であり、この性質が特殊化写像を特徴づける。
$Z' = f(Z) \subset Y$ とおく。このとき $Z'$ は $R$ を支配する $Y$ の
整閉部分スキームである。したがって $sp_{Y/R}([Z'_K]) = [Z'_\kappa]_r$ である。

\medskip\noindent
$[Z]$ を $Z_{r + 1}(X)$ の元とみなせる。定義により、
$\dim(Z') < r + 1$ ならば $f_*[Z] = 0$ であり、$Z \to Z'$ が次数 $d$ の
生成有限射ならば $f_*[Z] = d[Z']$ である。固有前送りは $Y_K \to Y$ による
平坦引き戻しと可換なので (Chow Homology, Lemma
\ref{chow-lemma-flat-pullback-proper-pushforward})、対応して
$f_{K, *}[Z_K] = 0$ または $f_{K, *}[Z_K] = d[Z'_K]$ となる。
可換図式
$$
\xymatrix{
X_\kappa \ar[d] \ar[r]_i & X \ar[d] \\
Y_\kappa \ar[r]^j & Y
}
$$
に Chow Homology, Lemma \ref{chow-lemma-closed-in-X-gysin} を適用する。
明らかに $i^*[Z] = [Z_k]_r$ および $j^*[Z'] = [Z'_\kappa]_r$ なので、
$f_{\kappa, *}[Z_\kappa]_r = 0$ または
$f_{\kappa, *}[Z_\kappa] = d[Z'_\kappa]_r$ を得る。以上を合わせれば結論が従う。
\end{proof}









\section{ファイバー上のサイクルの族}
\label{relative-cycles-section-cycles-fibres}

\noindent
$f : X \to S$ を局所有限型なスキームの射とし、$r \geq 0$ を整数とする。
{\it $X/S$ のファイバー上の $r$-サイクルの族 $\alpha$} とは、族
$$
\alpha = (\alpha_s)_{s \in S}
$$
であって、スキーム $S$ の点 $s$ により添字づけられ、
$\alpha_s \in Z_r(X_s)$ が $s$ における $f$ のスキーム論的ファイバー
$X_s$ 上の $r$-サイクルであるものをいう。ファイバー上の $r$-サイクルの族には、
いくつかの構成を施すことができる。

\medskip\noindent
{\bf 基底変換。} 図式
$$
\xymatrix{
X' \ar[r] \ar[d] & X \ar[d]^f \\
S' \ar[r]^g & S
}
$$
を、$f$ が局所有限型であるスキームの射のデカルト正方形とする。
$r \geq 0$ を整数とする。$X/S$ のファイバー上の $r$-サイクルの族 $\alpha$ に対し、
$\alpha$ の {\it 基底変換} $g^*\alpha$ を族
$$
g^*\alpha = (\alpha'_{s'})_{s' \in S'}
$$
と定義する。ここで $s = g(s')$ とおくと、$\alpha'_{s'} \in Z_r(X'_{s'})$ は
Section \ref{relative-cycles-section-relative-fields} におけるサイクル $\alpha_s$ の基底変換であり、
スキーム論的ファイバーの同一視
$X'_{s'} = X_s \times_{\Spec(\kappa(s))} \Spec(\kappa(s'))$
を用いる。

\medskip\noindent
{\bf 制限。} $f : X \to S$ を局所有限型なスキームの射とし、$r \geq 0$ を整数とする。
$U \subset X$ および $V \subset S$ を、$f(U) \subset V$ を満たす開部分スキームとする。
$X/S$ のファイバー上の $r$-サイクルの族 $\alpha$ に対し、$\alpha$ の
{\it 制限} $\alpha|_U$ を、$U/V$ のファイバー上の $r$-サイクルの族
$$
\alpha|_U = (\alpha_s|_{U_s})_{s \in V}
$$
と定義する。これはスキーム論的ファイバーへの制限からなる。

\medskip\noindent
{\bf 平坦引き戻し。} $X \to S$ を局所有限型なスキームの射とし、
$r, e \geq 0$ を整数とする。$f : X' \to X$ を、局所有限型かつ相対次元 $e$ の
平坦射とする。$X/S$ のファイバー上の $r$-サイクルの族 $\alpha$ に対し、
$\alpha$ の {\it 平坦引き戻し} $f^*\alpha$ を、ファイバー上の
$(r + e)$-サイクルの族
$$
f^*\alpha = (f_s^*\alpha_s)_{s \in S}
$$
と定義する。ここで $f_s^*\alpha_s \in Z_{r + e}(X'_s)$ は、スキーム論的
ファイバーの相対次元 $e$ の平坦射 $f_s : X'_s \to X_s$ による、
$Z_r(X_s)$ のサイクル $\alpha_s$ の平坦引き戻しである。

\medskip\noindent
{\bf 固有前送り。} 可換図式
$$
\xymatrix{
X \ar[rr]_f \ar[rd] & & Y \ar[ld] \\
& S
}
$$
において、$X$ と $Y$ が $S$ 上局所有限型であり、$f$ が固有であるとする。
$r \geq 0$ を整数とする。$X/S$ のファイバー上の $r$-サイクルの族 $\alpha$ に対し、
$\alpha$ の {\it 固有前送り} $f_*\alpha$ を、$Y/S$ のファイバー上の
$r$-サイクルの族
$$
f_*\alpha = (f_{s, *}\alpha_s)_{s \in S}
$$
と定義する。ここで $f_{s, *}\alpha_s \in Z_r(Y_s)$ は、スキーム論的ファイバーの
固有射 $f_s : X_s \to Y_s$ による、$Z_r(X_s)$ のサイクル $\alpha_s$ の
固有前送りである。

\begin{lemma}
\label{relative-cycles-lemma-compatibilities}
上の演算の間には、次の両立性がある：
(1) 基底変換は関手的である；
(2) 制限は、基底変換と平坦引き戻し（の特別な場合）の組合せである；
(3) 平坦引き戻しは基底変換と可換である；
(4) 平坦引き戻しは関手的である；
(5) 固有前送りは基底変換と可換である；
(6) 固有前送りは関手的である；
(7) 固有前送りは平坦引き戻しと可換である。
\end{lemma}

\begin{proof}
これらの両立性はそれぞれ、Chow homology の章で証明された対応する結果を、
考えているスキームの $S$ 上のファイバーに適用することで直ちに従う。
正確な主張と詳細な証明は省略する。いくつかの参照先を挙げる。
(1): Chow Homology, Lemma
\ref{chow-lemma-compose-base-change}.
(2): 明らかである。
(3): Chow Homology, Lemma
\ref{chow-lemma-pullback-base-change-pullback}.
(4): Chow Homology, Lemma
\ref{chow-lemma-compose-flat-pullback}.
(5): Chow Homology, Lemma
\ref{chow-lemma-pullback-base-change-pushforward}.
(6): Chow Homology, Lemma
\ref{chow-lemma-compose-pushforward}.
(7): Chow Homology, Lemma
\ref{chow-lemma-flat-pullback-proper-pushforward}.
\end{proof}

\begin{example}
\label{relative-cycles-example-family-associated-module}
$f : X \to S$ を局所有限型なスキームの射とし、$r \geq 0$ を整数とする。
$\mathcal{F}$ を有限型な準連接 $\mathcal{O}_X$-加群とする。$s \in S$ に対し、
$\mathcal{F}$ の $X_s$ への引き戻しを $\mathcal{F}_s$ と書く。
すべての $s \in S$ について $\dim(\text{Supp}(\mathcal{F}_s)) \leq r$ と仮定する。
このとき $\mathcal{F}$ に、次の式で定義される $X/S$ のファイバー上の
$r$-サイクルの族 $[\mathcal{F}/X/S]_r$ を対応させることができる：
$$
[\mathcal{F}/X/S]_r = ([\mathcal{F}_s]_r)_{s \in S}
$$
ここで $[\mathcal{F}_s]_r$ は Chow Homology, Definition
\ref{chow-definition-cycle-associated-to-coherent-sheaf} により与えられる。
\end{example}

\begin{lemma}
\label{relative-cycles-lemma-family-associated-module}
Example \ref{relative-cycles-example-family-associated-module} の構成は、基底変換、制限、
および平坦引き戻しと両立する。
\end{lemma}

\begin{proof}
Chow Homology, Lemmas
\ref{chow-lemma-pullback-coherent-base-change} および
\ref{chow-lemma-pullback-coherent} を参照せよ。
\end{proof}

\begin{example}
\label{relative-cycles-example-family-associated-closed}
$f : X \to S$ を局所有限型なスキームの射とし、$r \geq 0$ を整数とする。
$Z \subset X$ を閉部分スキームとする。$s \in S$ に対し、$X_s$ における $Z$ の
逆像を $Z_s$ と書く。これは同じことだが、$X_s$ の閉部分スキームとみなした
$s$ における $Z$ のスキーム論的ファイバーである。
すべての $s \in S$ について $\dim(Z_s) \leq r$ と仮定する。このとき $Z$ に、
次の式で定義される $X/S$ のファイバー上の $r$-サイクルの族 $[Z/X/S]_r$
を対応させることができる：
$$
[Z/X/S]_r = ([Z_s]_r)_{s \in S}
$$
ここで $[Z_s]_r$ は Chow Homology, Definition
\ref{chow-definition-cycle-associated-to-closed-subscheme} により与えられる。
\end{example}

\begin{lemma}
\label{relative-cycles-lemma-family-associated-closed}
Example \ref{relative-cycles-example-family-associated-closed} の構成は、基底変換、制限、
および平坦引き戻しと両立する。
\end{lemma}

\begin{proof}
$\mathcal{F} = (Z \to X)_*\mathcal{O}_Z$ とおけば、これは
Lemma \ref{relative-cycles-lemma-family-associated-module} の特別な場合である。Chow Homology,
Lemma \ref{chow-lemma-cycle-closed-coherent} も参照せよ。
\end{proof}

\begin{remark}[台]
\label{relative-cycles-remark-supports-family}
$f : X \to S$ を局所有限型なスキームの射とし、$r \geq 0$ を整数とする。
$\alpha$ を $X/S$ のファイバー上の $r$-サイクルの族とする。$\alpha$ の
{\it 台} を
$$
\text{Supp}(\alpha) =
\bigcup\nolimits_{s \in S} \text{Supp}(\alpha_s) \subset X
$$
と定義する。ここで $\text{Supp}(\alpha_s) \subset X_s$ はサイクル $\alpha_s$ の
台である (Chow Homology, Definition \ref{chow-definition-support-cycle} を参照せよ)。
台 $\text{Supp}(\alpha)$ が $X$ の閉部分集合であることは稀である。
\end{remark}

\begin{lemma}
\label{relative-cycles-lemma-support-family}
Remark \ref{relative-cycles-remark-supports-family} のように台を取る操作は、基底変換、制限、
および平坦引き戻しと両立する。
\end{lemma}

\begin{proof}
省略する。
\end{proof}

\begin{lemma}
\label{relative-cycles-lemma-coequalizer-dim-r}
$f : X \to S$ を局所有限型なスキームの射とし、$r \geq 0$ を整数とする。
$g : S' \to S$ をスキームの全射とする。$S'' = S' \times_S S'$ とおき、
$f' : X' \to S'$ および $f'' : X'' \to S''$ を $f$ の基底変換とする。
$x \in X$ が $\text{trdeg}_{\kappa(f(x))}(\kappa(x)) = r$ を満たすとする。
\begin{enumerate}
\item $x$ に写る $x' \in X'$ で
$\text{trdeg}_{\kappa(f'(x'))}(\kappa(x')) = r$ を満たすものが存在する。
\item $x'_1, x'_2 \in X'$ がともに (1) のような点ならば、
$\text{trdeg}_{\kappa(f''(x''))}(\kappa(x'')) = r$ および
$\text{pr}_i(x'') = x'_i$ を満たす $x'' \in X''$ が存在する。
\end{enumerate}
\end{lemma}

\begin{proof}
(1) は Morphisms, Lemma
\ref{morphisms-lemma-dimension-fibre-after-base-change} である。
$x'_1, x'_2$ を (2) のような点とする。$X'' = X' \times_X X'$ なので、
$x'_1$ と $x'_2$ の両方に写る $x'' \in X''$ が存在する
(例えば Descent, Lemma \ref{descent-lemma-equiv-fibre-product} を参照せよ)。
$x''$, $x'_i$, $x$ の像をそれぞれ $s'' \in S''$, $s'_i \in S'$, $s \in S$
と書く。$k = \kappa(s)$ とし、$Z \subset X_k$ を一般点が $x$ である
整閉部分スキームとする。このとき $x'_i$ は
$Z_{\kappa(s'_i)}$ のある既約成分の一般点である。
$Z'' \subset Z_{\kappa(s'')}$ を $x''$ を含む既約成分とし、その一般点を
$\xi'' \in Z''$ と書く。$\xi'' \leadsto x''$ なので、$\xi''$ は二つの射影の下で
$x'_i$ に写らなければならない。一方、この点は $Z$ の基底変換の既約成分の
一般点なので $\text{trdeg}_{\kappa(s'')}(\kappa(\xi'')) = r$ である。
\end{proof}

\begin{lemma}
\label{relative-cycles-lemma-descend-family}
$f : X \to S$ を局所有限型なスキームの射とし、$r \geq 0$ を整数とする。
$g : S' \to S$ をスキームの射とし、$X' = S' \times_S X$ とおく。
すべての $s \in S$ に対し、$g(s') = s$ かつ $\kappa(s')/\kappa(s)$ が
体の分離拡大となる点 $s' \in S'$ が存在すると仮定する。このとき：
\begin{enumerate}
\item $\alpha_1$ と $\alpha_2$ が $X/S$ のファイバー上の $r$-サイクルの族で、
$g^*\alpha_1 = g^*\alpha_2$ ならば $\alpha_1 = \alpha_2$ である。
\item $\alpha'$ を $X'/S'$ のファイバー上の $r$-サイクルの族とし、
$(S' \times_S S') \times_S X / (S' \times_S S')$ のファイバー上の
$r$-サイクルの族として
$\text{pr}_1^*\alpha' = \text{pr}_2^*\alpha'$ であるとする。このとき
$g^*\alpha = \alpha'$ を満たす $X/S$ のファイバー上の $r$-サイクルの族
$\alpha$ が一意に存在する。
\end{enumerate}
\end{lemma}

\begin{proof}
(1) は Section \ref{relative-cycles-section-relative-fields} で論じた基底変換写像の単射性から従う。
(この議論は $S' \to S$ が全射である限り成り立つ。)

\medskip\noindent
$\alpha'$ を (2) のような族とし、共通の値を
$\alpha'' = \text{pr}_1^*\alpha' = \text{pr}_2^*\alpha'$ と書く。

\medskip\noindent
$(X/S)^{(r)}$ を
$\text{trdeg}_{\kappa(f(x))}(\kappa(x)) = r$ を満たす $x \in X$ の集合とし、
$(X'/S')^{(r)}$ および $(X''/S'')^{(r)}$ も同様に定義する。係数を取ることで、
$\alpha'$ と $\alpha''$ を関数
$\alpha' : (X'/S')^{(r)} \to \mathbf{Z}$ および
$\alpha'' : (X''/S'')^{(r)} \to \mathbf{Z}$ とみなせる。関数
$$
\varphi : (X/S)^{(r)} \to \mathbf{Z}
$$
に対し、基底変換の操作にならって
$g^*\varphi : (X'/S')^{(r)} \to \mathbf{Z}$ を定義する。すなわち、
$x' \in (X'/S')^{(r)}$ が $x \in X$, $s' \in S'$, $s \in Z$ に写るとする。
一般点がそれぞれ $x'$ と $x$ である整閉部分スキームを
$Z' \subset X'_{s'}$ および $Z \subset X_s$ と書く。$\dim(Z') = r$ である。
$\dim(Z) < r$ ならば $(g^*\varphi)(x') = 0$ とおく。
$\dim(Z) = r$ ならば $Z'$ は $Z_{s'}$ の既約成分なので重複度
$m_{Z', Z_{s'}}$ をもつ。これを $m(x', g)$ と書き、
$$
(g^*\varphi)(x') = m(x', g) \varphi(x)
$$
と定義する。係数 $m(x', g)$ は常に正の整数であることに注意せよ
(例えば Lemma \ref{relative-cycles-lemma-multiplicities-field-extension} を参照せよ)。
同様に基底変換写像
$$
\text{pr}_1^*, \text{pr}_2^* :
\text{Map}((X'/S')^{(r)}, \mathbf{Z})
\longrightarrow
\text{Map}((X''/S'')^{(r)}, \mathbf{Z})
$$
を得る。基底変換の結合性から
$\text{pr}_1^* \circ g^* = \text{pr}_2^* \circ g^*$ が従う
(細部は省略する)。集合の写像で明示的に書けば、この等式は
$x'' \in (X''/S'')^{(r)}$ に対して
$$
m(x'', \text{pr}_1) m(\text{pr}_1(x''), g) =
m(x'', \text{pr}_2) m(\text{pr}_2(x''), g)
$$
が成り立つことを意味する。ただし $\text{pr}_1(x'')$ と
$\text{pr}_2(x'')$ は $(X''/S'')^{(r)}$ に属するとする。
Lemma \ref{relative-cycles-lemma-coequalizer-dim-r} と、上の表示式を用いる初等的な議論
\footnote{$x \in (X/S)^{(r)}$ に対し、$x$ に写る
$x' \in (X'/S')^{(r)}$ を選び、
$\alpha(x) = \alpha'(x')/m(x', g)$ とおく。この式と補題により良定義である。}
から、一意な写像
$$
\alpha : (X/S)^{(r)} \to \mathbf{Q}
$$
で $g^*\alpha = \alpha'$ を満たすものが存在する。証明を完了するには、
$\alpha$ が整数値を取ることを示せば十分である (細部を一つ省略した：
$\alpha$ が各ファイバー上で局所有限な和を定めることを確認する必要があるが、
これは $\alpha'$ についての対応する事実から従う)。像が $s \in S$ である任意の
$x \in (X/S)^{(r)}$ に対し、$\kappa(s')/\kappa(s)$ が分離的となる点
$s' \in S'$ を選べる。すると $s$ と $x$ に写る
$x' \in (X'/S')^{(r)}$ を選べ、この場合 $Z_{s'}$ は被約なので
$m(x', g) = 1$ である。したがって $\alpha(x) = \alpha'(x')$ は整数である。
\end{proof}

\begin{lemma}
\label{relative-cycles-lemma-pullback-universally-bijective}
$g : S' \to S$ を、剰余体の同型を誘導するスキームの全単射とする。
$f : X \to S$ を局所有限型とし、$X' = S' \times_S X$ とおく。
$r \geq 0$ とする。このとき $g$ による基底変換は、$X/S$ のファイバー上の
$r$-サイクルの族の群と $X'/S'$ のファイバー上の $r$-サイクルの族の群との間の
全単射を定める。
\end{lemma}

\begin{proof}
省略する。
\end{proof}








\section{相対サイクル}
\label{relative-cycles-section-families-specialization}

\noindent
以下の定義を用いる。\cite{SV} の定義との比較については
Section \ref{relative-cycles-section-compare} を参照せよ。

\begin{definition}
\label{relative-cycles-definition-relative-cycles}
$S$ を局所 Noether スキームとし、$f : X \to S$ を局所有限型なスキームの射とする。
$r \geq 0$ を整数とする。{\it $X/S$ 上の相対 $r$-サイクル} とは、
$X/S$ のファイバー上の $r$-サイクルの族 $\alpha$ であって、
$S'$ が離散付値環のスペクトルである任意の射 $g : S' \to S$ に対して
$$
sp_{X'/S'}(\alpha_\eta) = \alpha_0
$$
を満たすものをいう。ここで $sp_{X'/S'}$ は Section
\ref{relative-cycles-section-specialization} のものであり、$\alpha_\eta$ (resp.\ $\alpha_0$) は
$\alpha$ の基底変換 $g^*\alpha$ の $S'$ の一般点 (resp.\ 閉点) における値である。
$X/S$ 上のすべての相対 $r$-サイクルの群を $z(X/S, r)$ と書く。
\end{definition}

\begin{lemma}
\label{relative-cycles-lemma-relative-cycle-functoriality}
$\alpha$ を Definition \ref{relative-cycles-definition-relative-cycles} の意味での
$X/S$ 上の相対 $r$-サイクルとする。このとき $\alpha$ の任意の制限、基底変換、
平坦引き戻し、または固有前送りは相対 $r$-サイクルである。
\end{lemma}

\begin{proof}
平坦引き戻しについては Lemma \ref{relative-cycles-lemma-specialization-flat-pullback} を用いる。
制限は平坦引き戻しの特別な場合である。基底変換については基底変換の推移性を用いる。
固有前送りについては Lemma \ref{relative-cycles-lemma-specialization-proper-pushforward} を用いる。
\end{proof}

\begin{lemma}
\label{relative-cycles-lemma-relative-cycles-h-descent}
$f : X \to S$ をスキームの射とする。$S$ は局所 Noether で、$f$ は局所有限型であると
仮定する。$r \geq 0$ を整数とし、$\alpha$ を $X/S$ のファイバー上の
$r$-サイクルの族とする。$\{g_i : S_i \to S\}$ を h-被覆とする
(More on Flatness, Definition \ref{flat-definition-h-covering})。
このとき $\alpha$ が相対 $r$-サイクルであることと、各基底変換
$g_i^*\alpha$ が相対 $r$-サイクルであることは同値である。
\end{lemma}

\begin{proof}
$\alpha$ が相対 $r$-サイクルならば、Lemma
\ref{relative-cycles-lemma-relative-cycle-functoriality} により各基底変換 $g_i^*\alpha$ も
相対 $r$-サイクルである。逆に、各 $g_i^*\alpha$ が相対 $r$-サイクルであると仮定する。
$S'$ が離散付値環のスペクトルである射 $g : S' \to S$ を取る。
$S$ を $S'$ に、$X$ を $X' = X \times_S S'$ に、$\alpha$ を
$\alpha' = g^*\alpha$ に置き換え、h-被覆の基底変換が h-被覆であること
(More on Flatness, Lemma \ref{flat-lemma-h}) を用いると、次の段落で扱う問題に帰着する。

\medskip\noindent
$S$ を、閉点 $0$ と一般点 $\eta$ をもつ離散付値環のスペクトルとする。
$sp_{X/S}(\alpha_\eta) = \alpha_0$ を示せばよい。h-被覆は定義により V-被覆なので、
ある $i$ と $S_i$ の点の特殊化 $s' \leadsto s$ が存在して
$g_i(s') = \eta$ および $g_i(s) = 0$ を満たす
(Topologies, Lemma \ref{topologies-lemma-refine-qcqs-V} を参照せよ)。
Properties, Lemma \ref{properties-lemma-locally-Noetherian-specialization-dvr}
により、離散付値環のスペクトル $S'$ からの射 $h : S' \to S_i$ で、
一般点 $\eta'$ を $s'$ に、閉点 $0'$ を $s$ に写すものを取れる。
$\alpha' = h^*g_i^*\alpha$ と書く。仮定により
$sp_{X'/S'}(\alpha'_{\eta'}) = \alpha'_{0'}$ である。
$g = g_i \circ h : S' \to S$ は離散付値環の拡大から誘導されるスキームの射なので、
Lemma \ref{relative-cycles-lemma-specialization-extension} により $sp_{X/S}$ と $sp_{X'/S'}$ は
ファイバー上の基底変換写像と両立する。Section \ref{relative-cycles-section-relative-fields} で
論じたように基底変換写像 $Z_r(X_0) \to Z_r(X'_{0'})$ は単射なので、
$sp_{X/S}(\alpha_\eta) = \alpha_0$ を得る。
\end{proof}

\begin{lemma}
\label{relative-cycles-lemma-families-specialization-fppf-descent}
$f : X \to S$ をスキームの射とする。$S$ は局所 Noether で、$f$ は局所有限型であると
仮定する。$r, e \geq 0$ を整数とし、$\alpha$ を $X/S$ のファイバー上の
$r$-サイクルの族とする。$\{f_i : X_i \to X\}$ を、局所有限型かつ相対次元 $e$ の
平坦射からなる共同全射な族とする。このとき $\alpha$ が相対 $r$-サイクルであることと、
各平坦引き戻し $f_i^*\alpha$ が相対 $r$-サイクルであることは同値である。
\end{lemma}

\begin{proof}
$\alpha$ が相対 $r$-サイクルならば、Lemma
\ref{relative-cycles-lemma-relative-cycle-functoriality} により各引き戻し $f_i^*\alpha$ も
相対 $r$-サイクルである。逆に、各 $f_i^*\alpha$ が相対 $r$-サイクルであると仮定する。
$S'$ が離散付値環のスペクトルである射 $g : S' \to S$ を取る。
$S$ を $S'$ に、$X$ を $X' = X \times_S S'$ に、$\alpha$ を
$\alpha' = g^*\alpha$ に置き換えると、次の段落で扱う問題に帰着する。

\medskip\noindent
$S$ を、閉点 $0$ と一般点 $\eta$ をもつ離散付値環のスペクトルとする。
$sp_{X/S}(\alpha_\eta) = \alpha_0$ を示せばよい。$S$ の閉点への $f_i$ の
基底変換を $f_{i, 0} : X_{i, 0} \to X_0$ と書く。$f_{i, \eta}$ についても同様である。
すると
$$
f_{i, 0}^*sp_{X/S}(\alpha_\eta) =
sp_{X_i/S}(f_{i, \eta}^*\alpha_\eta) = f_{i, 0}^*\alpha_0
$$
である。実際、最初の等号は Lemma \ref{relative-cycles-lemma-specialization-flat-pullback} により、
二つ目は仮定により成り立つ。写像族
$f_{i, 0}^* : Z_r(X_0) \to Z_r(X_{i, 0})$ は共同単射である
($f_{i, 0}$ が共同全射であることによる) から、所望の結論を得る。
\end{proof}

\begin{lemma}
\label{relative-cycles-lemma-check-after-closed}
$S$ を局所 Noether スキームとし、$i : X \to Y$ を $S$ 上局所有限型なスキームの
閉埋め込みとする。$r \geq 0$ とし、$\alpha$ を $X/S$ のファイバー上の
$r$-サイクルの族とする。このとき $\alpha$ が $X/S$ 上の相対 $r$-サイクルであることと、
$i_*\alpha$ が $Y/S$ 上の相対 $r$-サイクルであることは同値である。
\end{lemma}

\begin{proof}
基底変換は $i_*$ と可換なので (Lemma \ref{relative-cycles-lemma-compatibilities})、次を示せば十分である：
$S$ が一般点 $\eta$ と閉点 $0$ をもつ離散付値環のスペクトルならば、
$sp_{X/S}(\alpha_\eta) = \alpha_0$ であることと
$sp_{Y/S}(i_{\eta, *}\alpha_\eta) = i_{0, *}\alpha_0$ であることは同値である。
これは、$i_{0, *} : Z_r(X_0) \to Z_r(Y_0)$ が単射であり、かつ Lemma
\ref{relative-cycles-lemma-specialization-proper-pushforward} により
$i_{0, *}sp_{X/S}(\alpha_\eta) = sp_{Y/S}(i_{\eta, *}\alpha_\eta)$
だからである。
\end{proof}

\noindent
次の補題は Lemma \ref{relative-cycles-lemma-relative-cycles-equal} で強められる。

\begin{lemma}
\label{relative-cycles-lemma-uniqueness-extension}
$f : X \to S$ をスキームの射とする。$S$ は局所 Noether で、$f$ は局所有限型であると
仮定する。$r \geq 0$ とし、$\alpha$ と $\beta$ を $X/S$ 上の相対
$r$-サイクルとする。次は同値である：
\begin{enumerate}
\item $\alpha = \beta$ である；
\item $S$ の既約成分の任意の一般点 $\eta \in S$ に対して
$\alpha_\eta = \beta_\eta$ である。
\end{enumerate}
\end{lemma}

\begin{proof}
(1) $\Rightarrow$ (2) は直ちに分かる。(2) を仮定する。各 $s \in S$ に対し、
$s$ へ特殊化する (2) のような $\eta$ を取れる。Properties, Lemma
\ref{properties-lemma-locally-Noetherian-specialization-dvr} により、
離散付値環のスペクトル $S'$ からの射 $g : S' \to S$ で、一般点 $\eta'$ を
$\eta$ に、閉点 $0$ を $s$ に写すものを取れる。このとき $\alpha_s$ と $\beta_s$ は
$Z_r(X_s)$ の元であり、$Z_r(X_{0'})$ の同じ元、すなわち
$sp_{X_{S'}/S'}(\alpha_{\eta'})$ に基底変換される。ここで $\alpha_{\eta'}$ は
$\alpha_\eta$ の基底変換である。Section \ref{relative-cycles-section-relative-fields} で論じたように
基底変換写像 $Z_r(X_s) \to Z_r(X_{0'})$ は単射なので、
$\alpha_s = \beta_s$ を得る。
\end{proof}

\begin{lemma}
\label{relative-cycles-lemma-family-associated-module-specialization}
Example \ref{relative-cycles-example-family-associated-module} の状況で、$S$ は局所 Noether であり、
$\mathcal{F}$ は次元 $\geq r$ において $S$ 上平坦であると仮定する
(More on Flatness, Definition \ref{flat-definition-flat-dimension-n})。
このとき $[\mathcal{F}/X/S]_r$ は $X/S$ 上の相対 $r$-サイクルである。
\end{lemma}

\begin{proof}
More on Flatness, Lemma \ref{flat-lemma-pre-flat-dimension-n} により、
$\mathcal{F}$ に関する仮定は任意の基底変換で保たれる。また Lemma
\ref{relative-cycles-lemma-family-associated-module} により、$[\mathcal{F}/X/S]_r$ の形成は
任意の基底変換と両立する。特殊化との両立性は離散付値環のスペクトルへの
基底変換後に検査されるので、$S$ が付値環のスペクトルである場合に帰着する。
この場合、集合
$U = \{x \in X \mid \mathcal{F}\text{ flat at }x\text{ over }S\}$
は More on Flatness, Lemma \ref{flat-lemma-finite-type-flat} により $X$ で開である。
仮定により $X$ における $U$ の補集合は $S$ 上で次元 $< r$ のファイバーをもつので、
包含 $U \subset X$ に沿う制限は、任意の基底変換後のファイバー上の
$r$-サイクルの群の同型を誘導し、これは特殊化写像および $\mathcal{F}$ に付随する
相対サイクルの形成と両立する。したがって
$[\mathcal{F}|_U / U /S]_r$ について特殊化との両立性を示せば十分である。
$\mathcal{F}|_U$ は $S$ 上平坦なので、これは Lemma
\ref{relative-cycles-lemma-specialization-module} と定義から従う。
\end{proof}

\begin{lemma}
\label{relative-cycles-lemma-family-associated-closed-specialization}
Example \ref{relative-cycles-example-family-associated-closed} の状況で、$S$ は局所 Noether であり、
$Z$ は次元 $\geq r$ において $S$ 上平坦であると仮定する。このとき
$[Z/X/S]_r$ は $X/S$ 上の相対 $r$-サイクルである。
\end{lemma}

\begin{proof}
仮定は $\mathcal{O}_Z$ が次元 $\geq r$ において $S$ 上平坦であることを意味する。
したがって $\mathcal{F} = (Z \to X)_*\mathcal{O}_Z$ として Lemma
\ref{relative-cycles-lemma-family-associated-module-specialization} を適用すれば結論を得る。
\end{proof}

\noindent
$S$ を局所 Noether スキームとし、$f : X \to S$ を有限型な射とする。
$r \geq 0$ とする。$Hilb(X/S, r)$ を、$Z \to S$ が平坦かつ相対次元
$\leq r$ である閉部分スキーム $Z \subset X$ の集合とする。Lemma
\ref{relative-cycles-lemma-family-associated-closed-specialization} により、各
$Z \in Hilb(X/S, r)$ に対して元 $[Z/X/S]_r \in z(X/S, r)$ を得る。
したがって群準同型
\begin{equation}
\label{relative-cycles-equation-cycle-classes}
\text{次を基底とする自由アーベル群： }Hilb(X/S, r) \longrightarrow z(X/S, r)
\end{equation}
を得る。これは $\sum n_i[Z_i]$ を $\sum n_i[Z_i/X/S]_r$ に送る。
相対 $r$-サイクルの重要な特徴は、適切な位相において $X$ と $S$ 上局所的に
この写像の像に入ることである。

\begin{lemma}
\label{relative-cycles-lemma-get-cycles}
$f : X \to S$ を、$S$ が Noether であるスキームの有限型射とする。
$r \geq 0$ とし、$\alpha$ を $X/S$ 上の相対 $r$-サイクルとする。このとき、
固有かつ完全分解的な
(More on Morphisms, Definition \ref{more-morphisms-definition-cd-morphism})
射 $g : S' \to S$ が存在し、$g^*\alpha$ は
(\ref{relative-cycles-equation-cycle-classes}) の像に入る。
\end{lemma}

\begin{proof}
Noether 帰納法により、同型でない任意の閉埋め込み $g : S' \to S$ による
$\alpha$ の引き戻しについて主張が成り立つと仮定してよい。

\medskip\noindent
$S_1 \subset S$ を既約成分とする (整閉部分スキームとみなす)。
$S_2 \subset S$ を $S'$ の補集合の閉包とする (被約閉部分スキームとみなす)。
$S_2 \not = \emptyset$ ならば、$S_1 \to S$ および $S_2 \to S$ による
$\alpha$ の引き戻しについて主張が成り立つ。対応する完全分解的な固有射を
$g_1 : S'_1 \to S_1$ および $g_2 : S'_2 \to S_2$ とすれば、
$S' = S'_1 \amalg S'_2 \to S$ は完全分解的な固有射であり、$S$ について主張が
成り立つことが分かる\footnote{$S'_1$ 上平坦かつ相対次元 $\leq r$ である
$S'_1 \times_S X$ の任意の閉部分スキームは、$S'$ 上平坦かつ相対次元
$\leq r$ である $S' \times_S X$ の閉部分スキームとみなせる。}。
したがって $S' \to S$ は全単射であると仮定でき、次の段落の場合に帰着する。

\medskip\noindent
$S$ が整であると仮定する。$\eta \in S$ を一般点とし、
$K = \kappa(\eta)$ を $S$ の関数体とする。このとき $\alpha_\eta$ は $X_K$ 上の
$r$-サイクルである。$\alpha_\eta = \sum n_i[Y_i]$ と書く。
$Y_i$ の閉包を取ると、$\eta$ への基底変換が $Y_i$ である整閉部分スキーム
$Z_i \subset X$ を得る。一般平坦性 (例えば Morphisms, Proposition
\ref{morphisms-proposition-generic-flatness}) により、各 $i$ について $Z_i$ は
$S$ の空でない開部分 $U$ 上平坦である。More on Flatness, Lemma
\ref{flat-lemma-flat-after-blowing-up} を適用すると、$Z_i$ の狭義変換
$Z'_i \subset X_{S'}$ が $S'$ 上平坦となるような $U$-許容ブローアップ
$g : S' \to S$ を取れる。このとき
$\beta = \sum n_i[Z'_i/X_{S'}/S']_r$ は
(\ref{relative-cycles-equation-cycle-classes}) の像に入り、Lemma
\ref{relative-cycles-lemma-uniqueness-extension} により $\beta = g^*\alpha$ である。

\medskip\noindent
しかし $S' \to S$ は完全分解的とは限らないので、これだけでは証明は完了しない。
これは容易に修正できる。$T \subset S$ を $U$ の補集合とし (閉部分スキームとみなす)、
Noether 帰納法により、$(T' \to S)^*\alpha$ が
(\ref{relative-cycles-equation-cycle-classes}) の像に入るような完全分解的な固有射
$T' \to T$ を取れる。このとき $S' \amalg T' \to S$ が所望の射である。
\end{proof}

\begin{lemma}
\label{relative-cycles-lemma-get-cycles-dvr}
$f : X \to S$ を有限型なスキームの射とし、$S$ は離散付値環のスペクトルであるとする。
$r \geq 0$ とする。このとき (\ref{relative-cycles-equation-cycle-classes}) は全射である。
\end{lemma}

\begin{proof}
これはもちろん Lemma \ref{relative-cycles-lemma-get-cycles} から従うが、次のように直接にも分かる。
$\alpha$ を $X/S$ 上の相対 $r$-サイクルとする。
$\alpha_\eta = \sum n_i[Z_i]$ と書く (和は有限である)。
Section \ref{relative-cycles-section-specialization} のように $Z_i$ の閉包を
$\overline{Z}_i \subset X$ と書く。このとき
$\alpha = \sum n_i[\overline{Z}_i/X/S]$.
\end{proof}

\begin{lemma}
\label{relative-cycles-lemma-relative-cycle-smooth}
$f : X \to S$ をスキームの射とし、$r \geq 0$ とする。$S$ は局所 Noether で、
$f$ は相対次元 $r$ の滑らかな射であると仮定する。
$\alpha \in z(X/S, r)$ とする。このとき $\alpha$ の台は $X$ で開かつ閉である
(より正確な結果については証明を参照せよ)。
\end{lemma}

\begin{proof}
$x \in X$ の像を $s \in S$ とする。$f$ は滑らかなので、$x$ を含む $X_s$ の
既約成分 $Z(x)$ は一意である。このとき $\dim(Z(x)) = r$ である。
サイクル $\alpha_s$ における $Z(x)$ の係数を $n_x$ とする。
関数 $x \mapsto n_x$ が $X$ 上局所定数であることを示す。

\medskip\noindent
$g : S' \to S$ を局所 Noether スキームの射とする。$X'$ を $X$ の基底変換、
$\alpha' = g^*\alpha$ を $\alpha$ の基底変換とする。
$x' \in X'$ が $s' \in S'$, $x \in X$, $s \in S$ に写るとする。
$n_{x'} = n_x$ であると主張する。実際、$Z(x)$ は $\kappa(s)$ 上滑らかなので
$Z(x) \times_{\Spec(\kappa(s))} \Spec(\kappa(s'))$
は被約である。$Z(x')$ はこのスキームの既約成分なので、Section
\ref{relative-cycles-section-relative-fields} の基底変換の定義により、$\alpha'_{s'}$ における
$Z(x')$ の係数 $n_{x'}$ は $\alpha_s$ における $Z(x)$ の係数 $n_x$ と同じである。
これで主張が示された。

\medskip\noindent
$X$ は局所 Noether なので、$x \mapsto n_x$ が局所定数であることを示すには、
$X$ における特殊化 $x' \leadsto x$ に対して $n_{x'} = n_x$ を示せば十分である。
$S'$ が離散付値環のスペクトルで、一般点 $\eta$ を $x'$ に、閉点 $0$ を $x$ に
写す射 $S' \to X$ を選ぶ。Properties, Lemma
\ref{properties-lemma-locally-Noetherian-specialization-dvr} を参照せよ。
すると $S' \to S$ による $f$ の基底変換 $X' \to S'$ は切断
$\sigma : S' \to X'$ をもち、$X'$ の点の特殊化
$\sigma(\eta) \leadsto \sigma(0)$ は $X$ の $x' \leadsto x$ に写る。
したがって次の段落の主張に帰着する。

\medskip\noindent
$S$ を一般点 $\eta$ と閉点 $0$ をもつ離散付値環のスペクトルとし、切断
$\sigma : S  \to X$ があるとする。$n_{\sigma(\eta)} = n_{\sigma(0)}$ と主張する。
More on Morphisms, Section \ref{more-morphisms-section-connected-components}
の議論、とりわけ More on Morphisms, Lemma
\ref{more-morphisms-lemma-connected-along-section-open} により、$X$ を開部分スキームで
置き換えれば、$X \to S$ のファイバーは連結であると仮定できる。これらのファイバーは
滑らかなので既約である。したがって $n = n_{\sigma(\eta)}$ とおくと
$\alpha_\eta = n[X_\eta]$ であり、関係 $sp_{X/S}(\alpha_\eta) = \alpha_0$ から
$\alpha_0 = n[X_0]$、すなわち所望の $n_{\sigma(0)} = n$ が従う。
\end{proof}

\begin{lemma}
\label{relative-cycles-lemma-relative-cycles-equal}
$f : X \to S$ をスキームの射とする。$S$ は局所 Noether で、$f$ は局所有限型であると
仮定する。$r \geq 0$ とし、$\alpha, \beta \in z(X/S, r)$ とする。
集合 $E = \{s \in S : \alpha_s = \beta_s\}$ は $S$ で閉である。
\end{lemma}

\begin{proof}
問題は $S$ 上局所的なので、$S$ はアフィンであると仮定してよい。
$X = \bigcup U_i$ をアフィン開被覆とする。
$E_i = \{s \in S : \alpha_s|_{U_{i, s}} = \beta_s|_{U_{i, s}}\}$ とおく。
このとき $E = \bigcap E_i$ である。したがって $U_i \to S$ と、$\alpha$, $\beta$ の
$U_i$ への制限について補題を示せば十分であり、次の段落の場合に帰着する。

\medskip\noindent
$X$ と $S$ が準コンパクトであると仮定する。$\gamma = \alpha - \beta$ とおくと
$E = \{s \in S : \gamma_s = 0\}$ である。Lemma
\ref{relative-cycles-lemma-family-associated-closed-specialization} により、$g_i^*\gamma$ が
(\ref{relative-cycles-equation-cycle-classes}) の像に入るような固有射の共同全射な有限族
$\{g_i : S_i \to S\}$ が存在する。$E_i = g_i^{-1}(E)$ は
$(g_i^*\gamma)_t = 0$ を満たす点 $t \in S_i$ の集合であることに注意せよ。
すべての $i$ について $E_i$ が閉ならば、$E = \bigcup g_i(E_i)$ も閉である。
したがって次の段落の場合に帰着する。

\medskip\noindent
$X$ と $S$ が準コンパクトであり、$S$ 上平坦かつ相対次元 $\leq r$ である有限個の
閉部分スキーム $Z_i \subset X$ によって
$\gamma = \sum n_i[Z_i/X/S]_r$ と書けると仮定する。
$X' = \bigcup Z_i$ とおく (スキーム論的合併)。このとき $i : X' \to X$ は閉埋め込みで、
$X'$ は $S$ 上相対次元 $\leq r$ である。また
$\gamma' = \sum n_i[Z_i/X'/S]_r$ とおけば $\gamma = i_*\gamma'$ である。
明らかに $E = E' = \{s \in S : \gamma'_s = 0\}$ なので、次の段落の場合に帰着する。

\medskip\noindent
$X$ は $S$ 上相対次元 $\leq r$ であると仮定する。
$s \in S$, $s \not \in E$ とする。$E \cap V$ が空となる $s$ の開近傍
$V \subset S$ が存在することを示す。$s \not \in E$ という仮定は、次元 $r$ の
整閉部分スキーム $Z \subset X_s$ で、$\gamma_s$ における $[Z]$ の係数 $n$ が
零でないものが存在することを意味する。$x \in Z$ を一般点とする。
$\dim(Z) = r$ なので、$x$ は $X_s$ のある既約成分 (すなわち $Z$) の一般点である。
したがって $X$ を $x$ の開近傍で置き換えれば、$Z$ が $X_s$ の唯一の既約成分であると
仮定できる。特に $\gamma_s = n[Z]$ である。

\medskip\noindent
ここで More on Morphisms, Lemma
\ref{more-morphisms-lemma-local-structure-finite-type} を適用すると、図式
$$
\xymatrix{
X \ar[dd] & X' \ar[l]^g \ar[d]^\pi & x \ar@{|->}[dd] &
x' \ar@{|->}[l]  \ar@{|->}[d] \\
& Y \ar[d]^h & & y \ar@{|->}[d] \\
S \ar@{=}[r] & S & s & s \ar@{=}[l]
}
$$
で、そこに列挙されたすべての性質をもつものを得る。
$\gamma' = g^*\gamma$ を平坦引き戻しとする。
$E \subset E' = \{s \in S: \gamma'_s = 0\}$ であり、$s \not \in E'$ であることに
注意せよ。実際、$Z' \subset X'_s$ を $x'$ の閉包とすると、$\gamma'_s$ における
$Z'$ の係数は零でない。同様に $\gamma'' = \pi_*\gamma'$ とおく。このとき
$E' \subset E'' = \{s \in S: \gamma''_s = 0\}$ であり、$s \not \in E''$ である。
実際、$Z'' \subset Y_s$ を $y$ の閉包とすると、$\gamma''_s$ における $Z''$ の係数は
零でない。Lemma \ref{relative-cycles-lemma-relative-cycle-smooth} と $Y \to S$ の開性により、
$s$ のある開近傍は $E''$ と交わらない。これで証明が完了する。
\end{proof}

\begin{lemma}
\label{relative-cycles-lemma-descend-through-limit}
$S = \lim_{i \in I} S_i$ を、遷移射がアフィンである Noether スキームの有向逆系の
極限とする。$0 \in I$ とし、$X_0 \to S_0$ をスキームの有限型射とする。
$i \geq 0$ に対し $X_i = S_i \times_{S_0} X_0$ とおき、
$X = S \times_{S_0} X_0$ とおく。$S$ も Noether ならば
$$
z(X/S, r) = \colim_{i \geq 0} z(X_i/S_i, r)
$$
である。ここで遷移写像は相対 $r$-サイクルの基底変換による。
\end{lemma}

\begin{proof}
$i \geq 0$ とし、$\alpha_i, \beta_i \in z(X_i/S_i, r)$ が
$z(X/S, r)$ の同じ元に写るとする。このとき $S \to S_i$ は Lemma
\ref{relative-cycles-lemma-relative-cycles-equal} の閉部分集合 $E \subset S_i$ の中に写る。
したがってある $j \geq i$ について射 $S_j \to S_i$ も $E$ の中に写る
(Limits, Lemma \ref{limits-lemma-limit-contained-in-constructible} を参照せよ)。
ゆえに $\alpha_i$ と $\beta_i$ の $S_j$ への基底変換は一致する。
したがって写像は単射である。

\medskip\noindent
$\alpha \in z(X/S, r)$ とする。Lemma \ref{relative-cycles-lemma-get-cycles} を適用すると、
$g^*\alpha$ が (\ref{relative-cycles-equation-cycle-classes}) の像に入るような完全分解的な固有射
$g : S' \to S$ を得る。$X' = S' \times_S X$ とおく。$S'$ 上平坦かつ相対次元
$\leq r$ である閉部分スキーム $Z_a \subset X'$ を用いて
$g^*\alpha = \sum n_a [Z_a/X'/S']_r$ と書く。

\medskip\noindent
ここで Limits, Section \ref{limits-section-descending-relative} の機構を用いる。
次を満たすある $i \geq 0$ を取れる：
\begin{enumerate}
\item $S$ への基底変換が $g : S' \to S$ である完全分解的な固有射
$g_i : S'_i \to S_i$；
\item $X'_i = S'_i \times_{S_i} X_i$ とおくと、$S'_i$ 上平坦かつ相対次元
$\leq r$ で、$S'$ への基底変換が $Z_a$ である閉部分スキーム
$Z_{ai} \subset X'_i$。
\end{enumerate}
このために Limits, Lemmas
\ref{limits-lemma-descend-finite-presentation},
\ref{limits-lemma-descend-closed-immersion-finite-presentation},
\ref{limits-lemma-descend-flat-finite-presentation},
\ref{limits-lemma-eventually-proper}, および
\ref{limits-lemma-limit-dimension}
と More on Morphisms, Lemma \ref{more-morphisms-lemma-descend-cd} を用いる。
次を考える：
$\alpha'_i = \sum n_a [Z_{ai}/X'_i/S'_i]_r \in z(X'_i/S'_i, r)$.
構成により $z(X'/S', r)$ における $\alpha'_i$ の像は基底変換 $g^*\alpha$ と一致する。

\medskip\noindent
$S''_i = S'_i \times_{S_i} S'_i$, $X''_i = S''_i \times_{S_i} X_i$ とおき、
$S'' = S' \times_S S'$, $X'' = S'' \times_S X$ とおく。射影を
$\text{pr}_1, \text{pr}_2 : S'' \to S'$ および
$\text{pr}_1, \text{pr}_2 : S''_i \to S'_i$ と書く。二つの基底変換
$\text{pr}_1^*\alpha'_i$ と $\text{pr}_1^*\alpha'_i$ は、
$\text{pr}_1^*g^*\alpha = \text{pr}_1^*g^*\alpha$ なので
$z(X''/S'', r)$ の同じ元に写る。したがって証明の最初の段落により、$i$ を大きくすれば
$\text{pr}_1^*\alpha'_i = \text{pr}_1^*\alpha'_i$ と仮定できる。
Lemma \ref{relative-cycles-lemma-descend-family} により、$g_i^*\alpha_i = \alpha'_i$ を満たす
$X_i/S_i$ のファイバー上の $r$-サイクルの族 $\alpha_i$ が一意に存在する
(ここで $S'_i \to S_i$ が完全分解的であることを用いる)。Lemma
\ref{relative-cycles-lemma-relative-cycles-h-descent} により
$\alpha_i \in z(X_i/S_i, r)$ である。Lemma \ref{relative-cycles-lemma-descend-family} の一意性から、
$z(X/S, r)$ における $\alpha_i$ の像は $\alpha$ である。これで証明が完了する。
\end{proof}

\begin{lemma}
\label{relative-cycles-lemma-thickening}
$S$ を局所 Noether スキームとし、$i : X \to X'$ を $S$ 上局所有限型なスキームの
厚化とする。$r \geq 0$ とする。このとき
$i_* : z(X/S, r) \to z(X'/S, r)$ は全単射である。
\end{lemma}

\begin{proof}
$i_s : X_s \to X'_s$ は厚化なので、$i_*$ が $X/S$ のファイバー上の
$r$-サイクルの族と $X'/S$ のファイバー上の $r$-サイクルの族との間の全単射を
誘導することは明らかである。また $\alpha$ を $X/S$ のファイバー上の
$r$-サイクルの族とすると、
$\alpha \in z(X/S, r) \Leftrightarrow i_*\alpha \in z(X'/S, r)$
が Lemma \ref{relative-cycles-lemma-check-after-closed} により成り立つ。補題が従う。
\end{proof}

\begin{lemma}
\label{relative-cycles-lemma-extend-to-larger}
$S$ を局所 Noether スキームとし、$X$ を $S$ 上局所有限型なスキームとする。
$r \geq 0$ とする。$U \subset X$ を、$X \setminus U$ が $S$ 上相対次元 $< r$、
すなわちすべての $s \in S$ について $\dim(X_s \setminus U_s) < r$ となる
開部分とする。このとき制限は全単射
$z(X/S, r) \to z(U/S, r)$ を定める。
\end{lemma}

\begin{proof}
次元に関する仮定により $Z_r(X_s) \to Z_r(U_s)$ は全単射なので、制限は
$X/S$ のファイバー上の $r$-サイクルの族と $U/S$ のファイバー上の
$r$-サイクルの族との間の全単射を誘導する。これらの制限写像
$Z_r(X_s) \to Z_r(U_s)$ は基底変換および特殊化と両立する
(Lemmas \ref{relative-cycles-lemma-compatibilities} および
\ref{relative-cycles-lemma-specialization-flat-pullback} を参照せよ)。これから補題は容易に従う。
細部は省略する。
\end{proof}

\begin{lemma}
\label{relative-cycles-lemma-seminormalize-base}
$g : S' \to S$ を、剰余体の同型を誘導する局所 Noether スキームの普遍同相写像とする。
$f : X \to S$ を局所有限型とし、$X' = S' \times_S X$ とおく。
$r \geq 0$ とする。このとき $g$ による基底変換は全単射
$z(X/S, r) \to z(X'/S', r)$ を定める。
\end{lemma}

\begin{proof}
Lemma \ref{relative-cycles-lemma-pullback-universally-bijective} により、$X/S$ のファイバー上の
$r$-サイクルの族の群と $X'/S'$ のファイバー上の $r$-サイクルの族の群との間に
全単射がある。$\alpha$ を $X/S$ のファイバー上の $r$-サイクルの族とし、
$\alpha' = g^*\alpha$ をその基底変換とする。$R$ が離散付値環ならば、任意の射
$h : \Spec(R) \to S$ は、一意な射 $h' : \Spec(R) \to S'$ によって
$g \circ h'$ と分解する。実際、射
$S' \times_S \Spec(R) \to \Spec(R)$ は剰余体上で全単射を誘導する普遍同相写像であり、
したがって切断をもつ (例えば $R$ が半正規環であることによる。Morphisms,
Section \ref{morphisms-section-seminormalization} を参照せよ)。ゆえに $\alpha$ が
特殊化と両立するという条件 (すなわち相対 $r$-サイクルであるという条件) は、
$\alpha'$ が特殊化と両立するという条件と同値である。
\end{proof}






\section{等次元相対サイクル}
\label{relative-cycles-section-equidimensional}

\noindent
定義は次のとおりである。

\begin{definition}
\label{relative-cycles-definition-equidimensional}
$f : X \to S$ をスキームの射とする。$S$ は局所 Noether で、$f$ は局所有限型であると
仮定する。$r \geq 0$ を整数とする。$X/S$ 上の相対 $r$-サイクル $\alpha$ の台
(Remark \ref{relative-cycles-remark-supports-family}) が、$S$ 上の相対次元が $\leq r$ である
閉部分集合 $W \subset X$ に含まれるとき、$\alpha$ は {\it 等次元的} であるという。
$X/S$ 上のすべての等次元相対 $r$-サイクルの群を $z_{equi}(X/S, r)$ と書く。
\end{definition}

\begin{example}
\label{relative-cycles-example-not-equidimensional}
\begin{reference}
\cite[Example 3.1.9]{SV}
\end{reference}
等次元的でない相対 $r$-サイクルが存在する。実際、$k$ を体とし、
$S = \Spec(k[x, y])$ 上の $X = \Spec(k[x, y, t])$ を考える。
$s$ を $S$ の点とし、$x$, $y$ の像を $a, b \in \kappa(s)$ と書く。
$X/S$ 上の $0$-サイクルの族 $\alpha$ を次で定義する：
\begin{enumerate}
\item $b = 0$ ならば $\alpha_s = 0$ とする；そうでなければ
\item $\alpha_s = [p] - [q]$ とする。ここで $p$ (resp.\ $q$) は
$t = a/b$ (resp.\ $t = (a + b^2)/b$) を満たす
$\Spec(\kappa(s)[t])$ の $\kappa(s)$-有理点である。
\end{enumerate}
これが特殊化と両立することの証明は読者に委ねる。その着想は、
$v(b) > 0$ を満たす $\kappa(s)$ 上の任意の付値 $v$ の剰余体上の
$\mathbf{P}^1$ において、$a/b$ と $(a + b^2)/b = a/b + b$ が同じ点へ
極限することである。一方、$\alpha$ の台の閉包は $(0, 0)$ 上のファイバー全体を含む。
\end{example}

\begin{lemma}
\label{relative-cycles-lemma-equidimensional-functoriality}
$f : X \to S$ をスキームの射とする。$S$ は局所 Noether で、$f$ は局所有限型であると
仮定する。$r \geq 0$ を整数とし、$\alpha$ を $X/S$ 上の相対 $r$-サイクルとする。
$\alpha$ が等次元的ならば、$\alpha$ の任意の制限、基底変換、または平坦引き戻しも
等次元的である。
\end{lemma}

\begin{proof}
省略する。
\end{proof}

\begin{lemma}
\label{relative-cycles-lemma-check-equidimensional}
$f : X \to S$ をスキームの射とする。$S$ は局所 Noether で、$f$ は局所有限型であると
仮定する。$r \geq 0$ を整数とし、$\alpha$ を $X/S$ 上の相対 $r$-サイクルとする。
$\alpha$ が等次元的であることを検査するには、$X$ と $S$ 上 Zariski 局所的に
調べればよい。
\end{lemma}

\begin{proof}
実際、$\alpha$ が等次元的であるという条件は、$\alpha$ の台の閉包が $S$ 上
相対次元 $\leq r$ であることにほかならない。閉包を取る操作は開部分への制限と
可換なので、補題が従う (細部は省略する)。
\end{proof}

\begin{lemma}
\label{relative-cycles-lemma-equidimensional-fppf-descent}
$f : X \to S$ をスキームの射とする。$S$ は局所 Noether で、$f$ は局所有限型であると
仮定する。$r \geq 0$ を整数とし、$\alpha$ を $X/S$ 上の相対 $r$-サイクルとする。
$\{g_i : S_i \to S\}$ を fppf 被覆とする。このとき $\alpha$ が等次元的であることと、
各基底変換 $g_i^*\alpha$ が等次元的であることは同値である。
\end{lemma}

\begin{proof}
$\alpha$ が等次元的ならば、Lemma \ref{relative-cycles-lemma-equidimensional-functoriality}
により各 $g_i^*\alpha$ も等次元的である。逆に各 $g_i^*\alpha$ が等次元的であると
仮定する。$X$ における $\text{Supp}(\alpha)$ の閉包を $W$ と書く。
$g_i : S_i \to S$ は平坦かつ局所有限表示なので普遍的に開であり、したがって射
$f_i : X_i = S_i \times_S X \to X$ も普遍的に開である。
$\alpha_i = g_i^*\alpha$ と書く。Lemma \ref{relative-cycles-lemma-support-family} により
$\text{Supp}(\alpha_i) = f_i^{-1}(\text{Supp}(\alpha))$ である。
$f_i$ は開なので、$W_i = f_i^{-1}(W)$ は $\text{Supp}(\alpha_i)$ の閉包である。
したがって仮定により射 $W_i \to S_i$ は相対次元 $\leq r$ である。Morphisms,
Lemma \ref{morphisms-lemma-dimension-fibre-after-base-change} と、射
$S_i \to S$ が共同全射であることから、$W \to S$ は相対次元 $\leq r$ である。
\end{proof}

\begin{lemma}
\label{relative-cycles-lemma-equidimensional-descent-pullbacks}
$f : X \to S$ をスキームの射とする。$S$ は局所 Noether で、$f$ は局所有限型であると
仮定する。$r, e \geq 0$ を整数とし、$\alpha$ を $X/S$ 上の相対
$r$-サイクルとする。$\{f_i : X_i \to X\}$ を、局所有限型かつ相対次元 $e$ の
平坦射からなる共同全射な族とする。このとき $\alpha$ が等次元的であることと、
各平坦引き戻し $f_i^*\alpha$ が等次元的であることは同値である。
\end{lemma}

\begin{proof}
省略する。ヒント：Lemma \ref{relative-cycles-lemma-equidimensional-fppf-descent} の証明と同様に、
$\alpha$ の台の閉包 $W$ の $f_i$ による逆像は、$f_i^*\alpha$ の台の閉包
$W_i$ であることを示す。すると、すべての $i$ について $W_i \to S$ が
相対次元 $\leq r + e$ ならば、$W \to S$ は相対次元 $\leq r$ である。
\end{proof}





\section{重み付けと相対零サイクル}
\label{relative-cycles-section-weightings}

\noindent
この節では、準有限射に対して相対零サイクルが More on Morphisms, Section
\ref{more-morphisms-section-weightings} で定義された重み付けと密接に関係することを示す。

\medskip\noindent
$S$ を局所 Noether スキームとし、$f : X \to S$ を局所準有限なスキームの射とする。
このとき $z(X/S, 0) = z_{equi}(X/S, 0)$ であり、$r > 0$ に対して
$z(X/S, r) = 0$ である。$\alpha \in z(X/S, 0)$ に対し、写像
$$
w_\alpha : X \longrightarrow \mathbf{Z},\quad
x \mapsto \alpha(x) [\kappa(x) : \kappa(s)]_i \quad\text{ただし }s = f(x)
$$
を定義する。ここで $\alpha(x)$ はファイバー $X_s$ 上の $0$-サイクル
$\alpha_s$ における $x$ の係数を表し、$[K : k]_i$ は体の有限拡大の非分離次数を表す。

\begin{lemma}
\label{relative-cycles-lemma-weightings-pre}
上の状況で、$g : S' \to S$ が局所 Noether スキームの射ならば
$w_{g^*\alpha} = w_\alpha \circ g'$ である。ここで
$g' : X' \to X$ は射影 $X' = S' \times_S X \to X$ である。
\end{lemma}

\begin{proof}
$x' \in X'$ の $S', S, X$ における像を $s', s, x$ とする。
$\kappa(s')/\kappa(s)$ による $[x]$ の基底変換における $[x']$ の係数は、
局所環 $(\kappa(s') \otimes_{\kappa(s)} \kappa(x))_\mathfrak q$ の長さである。
ここで $\mathfrak q$ は $x'$ に対応する素イデアルである。したがって
$$
[\kappa(x) : \kappa(s)]_i =
\text{length}((\kappa(s') \otimes_{\kappa(s)} \kappa(x))_\mathfrak q)
[\kappa(x') : \kappa(s')]_i
$$
$k/\kappa(s')$ を代数閉包とする。$\mathfrak q$ の上にある素イデアル
$\mathfrak p \subset k \otimes_{\kappa(s)} \kappa(x)$ を選ぶ。次を示せたとする：
$$
[\kappa(x) : \kappa(s)]_i =
\text{length}((k \otimes_{\kappa(s)} \kappa(x))_\mathfrak p)
\quad\text{and}\quad
[\kappa(x') : \kappa(s')]_i =
\text{length}((k \otimes_{\kappa(s')} \kappa(x'))_\mathfrak p)
$$
すると
$$
\text{length}((\kappa(s') \otimes_{\kappa(s)} \kappa(x))_\mathfrak q)
\text{length}((k \otimes_{\kappa(s')} \kappa(x'))_\mathfrak p)
=
\text{length}((k \otimes_{\kappa(s)} \kappa(x))_\mathfrak p)
$$
が Algebra, Lemma \ref{algebra-lemma-pullback-module} と
$\kappa(s') \otimes_{\kappa(s)} \kappa(x) \to k \otimes_{\kappa(s)} \kappa(x)$
の平坦性により成り立つので、主張を得る。二つの等式を示すには最初のものを
証明すれば十分である。$\kappa(x)/\kappa/\kappa(s)$ を Fields, Lemma
\ref{fields-lemma-separable-first} で構成された部分体とする。このとき
$$
k \otimes_{\kappa(s)} \kappa(x) =
\prod\nolimits_{\sigma : \kappa \to k}
k \otimes_{\sigma, \kappa} \kappa(x)
$$
であり、各因子は次数
$[\kappa(x) : \kappa] = [\kappa(x) : \kappa(s)]_i$ の局所環である。
これは所望の主張である。
\end{proof}

\begin{lemma}
\label{relative-cycles-lemma-weightings}
$S$ を局所 Noether スキームとし、$f : X \to S$ を局所準有限なスキームの射とする。
\begin{enumerate}
\item $\alpha \in z(X/S, 0)$ に対し、上で構成した写像
$w_\alpha : X \to \mathbf{Z}$ は重み付けである。
\item $X$ が準コンパクトならば、任意の重み付け $w : X \to \mathbf{Z}$ に対し、
ある $\alpha \in z(X/S, 0)$ について $nw = w_\alpha$ となる整数 $n > 0$ が存在する。
\item $S$ が $\mathbf{F}_p$ 上のスキームならば、(2) の整数 $n$ を素数 $p$ の
冪として選べる。
\end{enumerate}
\end{lemma}

\begin{proof}
$\alpha \in z(X/S, 0)$ とし、More on Morphisms, Definition
\ref{more-morphisms-definition-weighting} のような図式
$$
\xymatrix{
X \ar[d]_f & U \ar[l]^h \ar[d]^\pi \\
Y & V \ar[l]_g
}
$$
を選ぶ。基底変換 $g^*\alpha$ の制限を $\beta \in z(U/V, 0)$ と書く。
基底変換との両立性 (Lemma \ref{relative-cycles-lemma-weightings-pre}) により
$w_\beta = w_\alpha \circ h$ であるから、$\int_\pi w_\beta$ が $V$ 上局所定数であることを
示せば十分である。次に
\begin{align*}
\left( \int_\pi w_\beta \right)(v)
& =
\sum\nolimits_{u \in U, \pi(u) = v} 
\beta(u) [\kappa(u) : \kappa(v)]_i [\kappa(u) : \kappa(v)]_s \\
& =
\sum\nolimits_{u \in U, \pi(u) = v} \beta(u)[\kappa(u) : \kappa(v)]
\end{align*}
最後の式は $\pi_*\beta \in z(V/V, 0)$ における $v$ の係数である。Lemma
\ref{relative-cycles-lemma-relative-cycle-smooth} により、この関数は $V$ 上局所定数である。

\medskip\noindent
逆に、$w : X \to S$ を重み付けとし、$X$ は準コンパクトであるとする。
十分割り切れる整数 $n$ を選ぶ。$\alpha$ を、$s \in S$ に対して
$$
\alpha_s =
\sum\nolimits_{f(x) = s} \frac{n w(x)}{[\kappa(x) : \kappa(s)]_i} [x]
$$
となる $X/S$ のファイバー上の $0$-サイクルの族とする。これは $X_s$ 上の
零サイクルである。$f$ のファイバーは普遍的に有界なので (Morphisms, Lemma
\ref{morphisms-lemma-locally-quasi-finite-qc-source-universally-bounded})、すべての
$s \in S$ に対して右辺が整数係数をもつように $n$ を選べる。$\alpha$ が相対
$0$-サイクルであることを示せば、補題の最後の主張も従う。そのためには
$\alpha$ が離散付値環に沿う特殊化と両立することを示す必要がある。
$w_\alpha$ の構成にも Lemma \ref{relative-cycles-lemma-weightings-pre} の証明にも、$\alpha$ が
相対サイクルであることは使われていない。したがって $w_\alpha = nw$ であり、
この等式は基底変換後も保たれる。また重み付けの基底変換も重み付けである
(More on Morphisms, Lemma \ref{more-morphisms-lemma-weighting-base-change} を参照せよ)。
ゆえに次の段落で扱う問題に帰着する。

\medskip\noindent
$S$ を一般点 $\eta$ と閉点 $0$ をもつ離散付値環のスペクトルとする。
$w : X \to S$ を重み付けとし、$X$ は $S$ 上準有限であるとする。$\alpha$ を、
前の段落で (適切な $n$ に対して) 構成した $X/S$ のファイバー上の
$0$-サイクルの族とする。$sp_{X/S}(\alpha_\eta) = \alpha_0$ を示せばよい。
$\beta \in z(X/S, 0)$ を、$\beta_\eta = \alpha_\eta$ および
$\beta_0 = sp_{X/S}(\alpha_\eta)$ を満たす $X/S$ 上の相対 $0$-サイクルとする。
このとき $w' = w_\beta - nw : X \to \mathbf{Z}$ は (上の結果により) 重み付けであり、
$X$ のうち $\eta$ に写る点では零である。More on Morphisms, Lemma
\ref{more-morphisms-lemma-weighting-specialization} により $w' = 0$ である。
これは $w_\beta = nw$ を意味し、したがって所望の $\alpha = \beta$ を得る。
\end{proof}

\begin{example}
\label{relative-cycles-example-n-must-be-positive-sometimes}
$p$ を素数とし、$k'/k$ を次数 $p^f$ の体の純非分離拡大とする。
$X = \Spec(k')$, $S = \Spec(k)$ とおく。$X$ の唯一の点 $x$ を $1$ に送る写像
$w : X \to \mathbf{Z}$ は $X \to S$ の重み付けである。一方、群 $z(X/S, 0)$ は
サイクル $\alpha = [x]$ により自由に生成される。したがってこの場合、
Lemma \ref{relative-cycles-lemma-weightings} で使える最小の整数 $n$ は $n = p^f$ である。
\end{example}

\begin{example}
\label{relative-cycles-example-Merkurjev}
上の議論を用いると、A.S.~Merkurjev による例 \cite[Example 3.5.10]{SV} を
「説明」できる。$p$ を素数とし、$k = \mathbf{F}_p(a, b)$ を $a$, $b$ による
$\mathbf{F}_p$ の純超越拡大とする。次をおく：
$$
A = k[x, y, z]/(a x^p + b y^p - z^p)
$$
これは $k$ 上有限型の正規整域である。スペクトル $S = \Spec(A)$ は、剰余体が $k$ の
$A$ の極大イデアル $(x, y, z)$ に対応する唯一の特異点 $s$ を除いて正則である。
$k' = k(a^{1/p}, b^{1/p})$, $B = k'[x, y]$ とおく。包含 $k \to k'$ を用い、
$x$ を $x$ に、$y$ を $y$ に、$z$ を $a^{1/p}x + b^{1/p}y$ に送る写像
$A \to B$ を考える。$X = \Spec(B)$ とおき、環写像 $A \to B$ に対応する射
$f : X \to S$ を考える。More on Morphisms, Lemma
\ref{more-morphisms-lemma-weighting-quasi-finite-Noetherian} により得られる $f$ の
重み付け $w : X \to \mathbf{Z}$ は、$B$ の分数体が $A$ 上純非分離なので定数値
$1$ をもつ！ $V = f^{-1}(U) \subset X$ とおく。奇跡的平坦性により
$V \to U$ は次数 $p$ の平坦射である。したがってサイクル
$\alpha = [V/V/U]_0 \in z(V/U, 0)$ は $w_\alpha = pw|_V$ を満たす。
しかし $s$ の上にある $X$ の唯一の点 $x$ の剰余体は $k$ 上次数 $p^2$ の $k'$ なので、
$\beta$ で $z(X/S, 0)$ に属し、$\alpha$ へ制限されるものは存在しえない。実際、重み付け
$pw - w_\beta$ は More on Morphisms, Lemma
\ref{more-morphisms-lemma-weighting-specialization} により零となるはずであり、
上の補題の証明中の式から
$$
\beta_s = \frac{pw(x)}{[\kappa(x) : \kappa(s)]_i} [x] =
\frac{1}{p}[x]
$$
が $Z_0(X_s)$ で成り立つことになるが、これは不可能である
\footnote{\cite{SV} では $pw$ に対応する相対サイクルの群 $Cycl(X/S, 0)$ の元が
実際に存在し、このサイクルは $X_s$ 上の非整数係数をもつサイクルへ特殊化する。}。
もちろんサイクル $p\alpha$ は $z(X/S, 0)$ の一意な元の制限である。
\end{example}









\section{有効相対サイクル}
\label{relative-cycles-section-effective}


\noindent
定義は次のとおりである。

\begin{definition}
\label{relative-cycles-definition-effective}
$f : X \to S$ をスキームの射とする。$S$ は局所 Noether で、$f$ は局所有限型であると
仮定する。$r \geq 0$ を整数とする。すべての $s \in S$ について $\alpha_s$ が
有効サイクル (Chow Homology, Definition \ref{chow-definition-effective-cycle})
であるとき、$X/S$ 上の相対 $r$-サイクル $\alpha$ は {\it 有効} であるという。
$X/S$ 上のすべての有効相対 $r$-サイクルのモノイドを $z^{eff}(X/S, r)$ と書く。
\end{definition}

\noindent
以下で、有効相対サイクルが等次元的であることを示す
(Lemma \ref{relative-cycles-lemma-effective-equidimensional} を参照せよ)。

\begin{lemma}
\label{relative-cycles-lemma-effective-functoriality}
$f : X \to S$ をスキームの射とする。$S$ は局所 Noether で、$f$ は局所有限型であると
仮定する。$r \geq 0$ を整数とし、$\alpha$ を $X/S$ 上の相対 $r$-サイクルとする。
$\alpha$ が有効ならば、$\alpha$ の任意の制限、基底変換、平坦引き戻し、または
固有前送りも有効である。
\end{lemma}

\begin{proof}
省略する。
\end{proof}

\begin{lemma}
\label{relative-cycles-lemma-check-effective}
$f : X \to S$ をスキームの射とする。$S$ は局所 Noether で、$f$ は局所有限型であると
仮定する。$r \geq 0$ を整数とし、$\alpha$ を $X/S$ 上の相対 $r$-サイクルとする。
$\alpha$ が有効であることを検査するには、$X$ と $S$ 上 Zariski 局所的に調べればよい。
\end{lemma}

\begin{proof}
省略する。
\end{proof}

\begin{lemma}
\label{relative-cycles-lemma-effective-descent}
$f : X \to S$ をスキームの射とする。$S$ は局所 Noether で、$f$ は局所有限型であると
仮定する。$r \geq 0$ を整数とし、$\alpha$ を $X/S$ 上の相対 $r$-サイクルとする。
$g : S' \to S$ を全射とする。このとき $\alpha$ が有効であることと、基底変換
$g^*\alpha$ が有効であることは同値である。
\end{lemma}

\begin{proof}
省略する。
\end{proof}

\begin{lemma}
\label{relative-cycles-lemma-effective-descent-pullbacks}
$f : X \to S$ をスキームの射とする。$S$ は局所 Noether で、$f$ は局所有限型であると
仮定する。$r, e \geq 0$ を整数とし、$\alpha$ を $X/S$ 上の相対
$r$-サイクルとする。$\{f_i : X_i \to X\}$ を、局所有限型かつ相対次元 $e$ の
平坦射からなる共同全射な族とする。このとき $\alpha$ が有効であることと、
各平坦引き戻し $f_i^*\alpha$ が有効であることは同値である。
\end{lemma}

\begin{proof}
省略する。
\end{proof}

\begin{lemma}
\label{relative-cycles-lemma-effective-support-closed}
$f : X \to S$ をスキームの射とする。$S$ は局所 Noether で、$f$ は局所有限型であると
仮定する。$r, e \geq 0$ を整数とし、$\alpha$ を $X/S$ 上の相対
$r$-サイクルとする。$\alpha$ が有効ならば、$\text{Supp}(\alpha)$ は $X$ で閉である。
\end{lemma}

\begin{proof}
$g : S' \to S$ を、整閉部分スキームとみなした既約成分の包含とする。
Lemmas \ref{relative-cycles-lemma-effective-functoriality} および \ref{relative-cycles-lemma-support-family} により、
基底変換 $g^*\alpha$ の台が $S' \times_S S$ で閉であることを示せば十分である。
したがって $S$ は一般点 $\eta$ をもつ整スキームであると仮定してよい。
$\text{Supp}(\alpha)$ が $\text{Supp}(\alpha_\eta)$ の閉包であることを示す。
そのため任意の $s \in S$ を取る。$S'$ が離散付値環のスペクトルで、一般点
$\eta' \in S'$ を $\eta$ に、閉点 $0 \in S'$ を $s$ に写す射 $g : S' \to S$ を取れる
(Properties, Lemma \ref{properties-lemma-locally-Noetherian-specialization-dvr}
を参照せよ)。すると $g^*\alpha$ の台が
$\text{Supp}((g^\alpha)_{\eta'})$ の閉包に等しいことを示せば十分であり、
次の段落の場合に帰着する。

\medskip\noindent
ここで $S$ は一般点 $\eta$ と閉点 $0$ をもつ離散付値環のスペクトルである。
$\text{Supp}(\alpha)$ が $\text{Supp}(\alpha_\eta)$ の閉包であることを示せばよい。
$\alpha$ は有効なので、$n_i > 0$ かつ $Z_i \subset X_\eta$ が次元 $r$ の
整閉部分スキームとなるように $\alpha_\eta = \sum n_i[Z_i]$ と書ける。
$\alpha_0 = sp_{X/S}(\alpha_\eta)$ なので、$\overline{Z}_i$ を $Z_i$ の閉包とすると
$\alpha_0 = \sum n_i [\overline{Z}_{i, 0}]_r$ である。Varieties, Lemma
\ref{varieties-lemma-dominate-valuation-ring-dimension-fibres} により
$\overline{Z}_{i, 0}$ は次元 $r$ の等次元スキームである。$n_i > 0$ なので
$\text{Supp}(\alpha_0)$ は $\overline{Z}_{i, 0}$ の合併に等しい。これは
$\bigcup \overline{Z}_i$ の $0$ 上のファイバーであり、さらに
$\bigcup Z_i$ の閉包である。これが所望の主張である。
\end{proof}

\begin{lemma}
\label{relative-cycles-lemma-effective-equidimensional}
$f : X \to S$ をスキームの射とする。$S$ は局所 Noether で、$f$ は局所有限型であると
仮定する。$r, e \geq 0$ を整数とし、$\alpha$ を $X/S$ 上の相対
$r$-サイクルとする。$\alpha$ が有効ならば $\alpha$ は等次元的である。
\end{lemma}

\begin{proof}
$\alpha$ が有効であると仮定する。Lemma \ref{relative-cycles-lemma-effective-support-closed} により
台 $\text{Supp}(\alpha)$ は $X$ で閉である。
$\text{Supp}(\alpha) \to S$ のファイバーはサイクル $\alpha_s$ の台であり、
したがって次元 $r$ なので、$\alpha$ は等次元的である。
\end{proof}

\begin{remark}
\label{relative-cycles-remark-representable}
$f : X \to S$ をスキームの射とし、$S$ は局所 Noether、$f$ は局所有限型であるとする。
反変関手
$$
\begin{matrix}
\text{局所有限型なスキーム }S'\text{ で} \\
\text{基底が }S
\end{matrix}
\longrightarrow
z^{eff}(X'/S', r)\text{、ただし }X' = S' \times_S X
$$
が表現可能かどうかを問うことができる。しかし
$z(X'/S', r) = z(X'_{red}/S'_{red}, r)$ なので、これは成り立ちえない
(実際の反例を作ることは読者に委ねる)。より適切な問いは、この関手が表現可能となる
左辺の部分圏を見つけられるかどうかであろう。Lemma
\ref{relative-cycles-lemma-seminormalize-base} は、少なくとも $S$ 上の半正規スキームの圏に
制限すべきことを示唆している。

\medskip\noindent
$S/\Spec(\mathbf{Q})$ が Nagata で $f$ が射影的ならば、
$S' \mapsto z^{eff}(X'/S', r)$ は半正規な $S'$ の圏上で表現可能であることが分かる。
大まかに言えば、これは \cite[Theorem 3.21]{KRC} の内容である。

\medskip\noindent
$S$ に正標数の点がある場合、半正規性を弱正規性で置き換えてもこれは成り立たない。
局所 Noether スキーム $T$ は、任意の双有理な普遍同相写像 $T' \to T$ が切断をもつとき
弱正規であるという。例として、標数 $2$ の体 $k$ に対し、
$S = \Spec(k)$ 上の $X = \mathbf{A}^2_k$ の次数 $2$ の $0$-サイクルを考える。
すなわち $W = X \times_S X$ 上には標準的な相対 $0$-サイクル
$\alpha \in z^{eff}(X_W/W, 0)$ があり、$w = (x_1, x_2) \in W = X^2$ に対して
$\alpha_w = [x_1] + [x_2]$ である。このサイクルは因子を交換する対合
$\sigma : W \to W$ の下で不変である。$W$ は滑らかであり (したがって正規、さらに
弱正規)、もし $z(-/-, r)$ が $k$ 上有限型な弱正規スキームの圏上で $M$ により
表現されるならば、$W$ から $M$ への $\sigma$-不変な射を得る。これはさらに商スキーム
$\text{Sym}^2_S(X) = W/\langle \sigma \rangle$ から $M$ への射を定める。
$\text{Sym}^2_S(X)$ は正規なので、$M$ のモジュライ性質により、$W$ への引き戻しが
$\alpha$ である
$X \times_S \text{Sym}^2_S(X) / \text{Sym}^2_S(X)$ 上の相対 $0$-サイクル
$\beta$ を得るはずである。しかし、そのようなサイクル $\beta$ は存在しない。
実際、$X = \Spec(k[u, v])$ と書くと、スキーム $\text{Sym}^2_S(X)$ は
$$
k[u_1 + u_2, u_1u_2, v_1 + v_2, v_1v_2, u_1v_1 + u_2v_2]
\subset
k[u_1, u_2, v_1, v_2]
$$
のスペクトルである。$\text{Sym}^2_S(X)$ における対角線
$u_1 = u_2, v_1 = v_2$ の像は閉部分スキーム
$V = \Spec(k[u_1^2, v_1^2])$ である。ここで $k$ の標数が $2$ であることを用いた。
$V$ の一般点 $\eta$ で見ると、サイクル $\beta_\eta$ は
$\mathbf{A}^2_{k(u_1^2, v_1^2)}$ 上の次数 $2$ の零サイクルであり、
$\mathbf{A}^2_{k(u_1, u_2)}$ への引き戻しは
$2[\text{点 }(u_1, v_2)]$ となるはずである。
これは明らかに不可能である。

\medskip\noindent
上の議論は \cite[Theorem 4.13]{KRC} と矛盾しない。その定理の Chow 多様体は関手を
粗く表現するだけだからである (実際には相異なる二つの関手であり、少し調べれば分かる
ように、そのうち一つだけが射影的な $X$ に対して我々の関手と一致する)。同様に
\cite[Section 4.4]{SV} では、射影的な $X/S$ に対して前層
$S' \mapsto z^{eff}(S' \times_S X/S', r)$ の $h$-層化が、表現可能関手の
$h$-層化に等しいことが示されている。
\end{remark}

\begin{remark}
\label{relative-cycles-remark-equidimensional-over-geometrically-unibranch}
$f : X \to S$ をスキームの射とし、$r \geq 0$ とする。
$Z \subset X$ を閉部分スキームとする。次を仮定する：
\begin{enumerate}
\item $S$ は Noether かつ幾何学的に単枝である；
\item $f$ は有限型である；
\item $Z \to S$ は相対次元 $\leq r$ である。
\end{enumerate}
このとき、十分割り切れるすべての整数 $n \geq 1$ に対し、$S$ の各一般点 $\eta$ で
$\alpha_\eta = n[Z_\eta]_r$ を満たす $X/S$ 上の有効相対 $r$-サイクル
$\alpha$ が一意に存在する。これは \cite[Theorem 3.4.2]{SV} の言い換えである。
この結果が必要になった場合には、ここで正確に定式化して証明する。
\end{remark}











\section{固有相対サイクル}
\label{relative-cycles-section-proper}

\noindent
我々の状況では、次がおそらく適切な定義である。

\begin{definition}
\label{relative-cycles-definition-proper}
$f : X \to S$ をスキームの射とする。$S$ は局所 Noether で、$f$ は局所有限型であると
仮定する。$r \geq 0$ を整数とする。$X/S$ 上の相対 $r$-サイクル $\alpha$ の台
(Remark \ref{relative-cycles-remark-supports-family}) が、$S$ 上固有な閉部分集合 $W \subset X$
(Cohomology of Schemes, Definition \ref{coherent-definition-proper-over-base})
に含まれるとき、$\alpha$ は {\it 固有相対サイクル} であるという。
$X/S$ 上のすべての固有相対 $r$-サイクルの群を $c(X/S, r)$ と書く。
\end{definition}

\noindent
Cohomology of Schemes, Lemma
\ref{coherent-lemma-closed-closed-proper-over-base} により、これは台の閉包が基底上
固有であることにほかならない。これらが群をなすことは Cohomology of Schemes,
Lemma \ref{coherent-lemma-union-closed-proper-over-base} から分かる。

\begin{lemma}
\label{relative-cycles-lemma-proper-functoriality}
$f : X \to S$ をスキームの射とする。$S$ は局所 Noether で、$f$ は局所有限型であると
仮定する。$r \geq 0$ を整数とし、$\alpha$ を $X/S$ 上の相対 $r$-サイクルとする。
$\alpha$ が固有ならば、$\alpha$ の任意の基底変換も固有である。
\end{lemma}

\begin{proof}
省略する。
\end{proof}

\begin{lemma}
\label{relative-cycles-lemma-proper-h-descent}
$f : X \to S$ をスキームの射とする。$S$ は局所 Noether で、$f$ は局所有限型であると
仮定する。$r \geq 0$ を整数とし、$\alpha$ を $X/S$ 上の相対 $r$-サイクルとする。
$\{g_i : S_i \to S\}$ を h-被覆とする。このとき $\alpha$ が固有であることと、
各基底変換 $g_i^*\alpha$ が固有であることは同値である。
\end{lemma}

\begin{proof}
$\alpha$ が固有ならば、Lemma \ref{relative-cycles-lemma-proper-functoriality} により各
$g_i^*\alpha$ も固有である。逆に各 $g_i^*\alpha$ が固有であると仮定する。
$\alpha$ が固有であることを示すには、明らかに $S$ 上アフィン局所的に調べれば十分である。
したがって $S$ はアフィンであると仮定する。このとき被覆 $\{S_i \to S\}$ を族
$\{T_j \to S\}$ で細分できる。ここで $g : T \to S$ は固有全射であり、
$T = \bigcup T_j$ は開被覆である。したがって $\beta = g^*\alpha$ は
$Y = T \times_S X$ 上、$T$ に対して固有である。Lemma
\ref{relative-cycles-lemma-support-family} により、$\beta$ の台は射 $f : Y \to X$ による
$\alpha$ の台の逆像である。したがって $f^{-1}\text{Supp}(\alpha)$ の閉包
$W \subset Y$ は $T$ 上固有である。射 $T \to S$ は固有なので $W$ は $S$ 上固有である。
すると Cohomology of Schemes, Lemma
\ref{coherent-lemma-functoriality-closed-proper-over-base} により、像
$f(W) \subset X$ は $S$ 上固有な閉部分集合である。$f(W)$ は
$\text{Supp}(\alpha)$ を含むので、$\alpha$ は固有である。
\end{proof}



\section{固有かつ等次元な相対サイクル}
\label{relative-cycles-section-proper-equidimensional}

\noindent
$f : X \to S$ をスキームの射とする。$S$ は局所 Noether で、$f$ は局所有限型であると
仮定する。$r \geq 0$ を整数とする。$X/S$ 上の相対 $r$-サイクル $\alpha$ が
等次元的 (Definition \ref{relative-cycles-definition-equidimensional}) かつ固有
(Definition \ref{relative-cycles-definition-proper}) であるとき、$\alpha$ は
{\it 固有かつ等次元な相対サイクル} であるという。$X/S$ 上のすべての固有かつ
等次元な相対 $r$-サイクルの群を $c_{equi}(X/S, r)$ と書く。

\medskip\noindent
同様に、$X/S$ 上の相対 $r$-サイクル $\alpha$ が有効
(Definition \ref{relative-cycles-definition-effective}) かつ固有
(Definition \ref{relative-cycles-definition-proper}) であるとき、$\alpha$ は
{\it 固有かつ有効な相対サイクル} であるという。$X/S$ 上のすべての固有かつ
有効な相対 $r$-サイクルのモノイドを $c^{eff}(X/S, r)$ と書く。Lemma
\ref{relative-cycles-lemma-effective-equidimensional} により、これらは等次元的であることに注意せよ。

\medskip\noindent
したがって包含写像の図式
$$
\xymatrix{
c^{eff}(X/S, r) \ar[r] \ar[d] &
c_{equi}(X/S, r) \ar[r] \ar[d] &
c(X/S, r) \ar[d] \\
z^{eff}(X/S, r) \ar[r] &
z_{equi}(X/S, r) \ar[r] &
z(X/S, r)
}
$$



\section{サイクルへの作用}
\label{relative-cycles-section-action}

\noindent
$S$ を、次元関数 $\delta$ を備えた局所 Noether かつ普遍鎖状なスキームとする
(Chow Homology, Section \ref{chow-section-setup} を参照せよ)。
$X \to Y$ を $S$ 上のスキームの射とし、その両端はともに $S$ 上局所有限型とする。
$r \geq 0$ とし、$\alpha$ を $X/Y$ のファイバー上の $r$-サイクルの族とする。
$e \in \mathbf{Z}$ に対して演算
$$
\alpha \cap - : Z_e(Y) \longrightarrow Z_{r + e}(X)
$$
を構成する。実際、$\beta \in Z_e(Y)$ に対し $\beta = \sum n_i[Z_i]$ と書く。
ここで $Z_i \subset Y$ は $\delta$-次元 $e$ の整閉部分スキームであり、族 $Z_i$ は
スキーム $Y$ で局所有限である。$y_i \in Z_i$ を一般点とし、
$\alpha_{y_i} = \sum m_{ij} [V_{ij}]$ と書く。したがって
$V_{ij} \subset X_{y_i}$ は次元 $r$ の整閉部分スキームであり、族 $V_{ij}$ は
スキーム $X_{y_i}$ で局所有限である。そこで
$$
\alpha \cap \beta = \sum n_i m_{ij} [\overline{V}_{ij}]
\quad\in\quad
Z_{r + e}(X)
$$
とおく。ここで $\overline{V}_{ij} \subset X$ は射
$V_{ij} \to X_{y_i} \to X$ のスキーム論的像である。同じことだが、
$\overline{V}_{ij} \subset X$ は $Z_i \subset Y$ を支配的に写し、一般ファイバーが
$V_{ij}$ である整閉部分スキームである。
$\dim_\delta(\overline{V}_{ij}) = r + e$ であり、閉部分スキームの族
$\overline{V}_{ij} \subset X$ が局所有限であることは容易に従う
(検証は省略する)。したがって $\alpha \cap \beta$ は確かに
$Z_{r + e}(X)$ の元である。

\begin{lemma}
\label{relative-cycles-lemma-action-bilinear}
上の構成は双線形である。すなわち
$(\alpha_1 + \alpha_2) \cap \beta = \alpha_1 \cap \beta +
\alpha_2 \cap \beta$ および $\alpha \cap (\beta_1 + \beta_2) =
\alpha \cap \beta_1 + \alpha \cap \beta_2$.
\end{lemma}

\begin{proof}
省略する。
\end{proof}

\begin{lemma}
\label{relative-cycles-lemma-action-opens}
$U \subset X$ と $V \subset Y$ が開で $f(U) \subset V$ ならば、
$(\alpha \cap \beta)|_U$ は $\alpha|_U \cap \beta|_V$ に等しい。
\end{lemma}

\begin{proof}
上で与えた $\alpha \cap \beta$ の明示的な記述から直ちに従う。
\end{proof}

\begin{lemma}
\label{relative-cycles-lemma-action-base-change}
$\alpha \cap \beta$ の形成は平坦基底変換および平坦引き戻しと両立する
(詳細については証明を参照せよ)。
\end{lemma}

\begin{proof}
$(S, \delta)$, $(S', \delta')$, $g : S' \to S$, $c \in \mathbf{Z}$ を
Chow Homology, Situation \ref{chow-situation-setup-base-change} のようなものとする。
$X \to Y$ を $S$ 上局所有限型なスキームの射とし、$X' \to Y'$ を
$X \to Y$ の $g$ による基底変換と書く。$\alpha$ を $X/Y$ のファイバー上の $r$-サイクルの族とし、
$\beta \in Z_e(Y)$ とする。$Y' \to Y$ による $\alpha$ の基底変換を $\alpha'$ と書き、
$g$ による $\beta$ の引き戻しを
$\beta' = g^*\beta \in Z_{e + c}(Y')$ と書く
(Chow Homology, Section \ref{chow-section-change-base} を参照せよ)。
基底変換との両立性とは、$\alpha' \cap \beta'$ が
$\alpha \cap \beta$ の基底変換であることを意味する。

\medskip\noindent
基底変換との両立性を証明する。$X'$ 上のサイクルの等式を証明するので、Lemma
\ref{relative-cycles-lemma-action-opens} により $Y$ 上局所的に調べてよい。したがって $Y$ は
アフィンであると仮定する。特に $\beta$ は素サイクルの有限線形結合である。
$- \cap -$ は第二変数について線形なので (Lemma \ref{relative-cycles-lemma-action-bilinear})、
$\delta$-次元 $e$ のある整閉部分スキーム $Z \subset Y$ に対し
$\beta = [Z]$ である場合に等式を証明すれば十分である。

\medskip\noindent
$y \in Z$ を一般点とする。$\alpha_y = \sum m_j [V_j]$ と書き、
$\overline{V}_j$ を $X$ における $V_j$ の閉包とする。このとき
$$
\alpha \cap \beta = \sum m_j[\overline{V}_j]
$$
$Y' = Y \times_S S'$ 上のサイクルとして、$\beta$ の基底変換は
$\beta' = \sum [Z \times_S S']_{e + c}$ である。
$Z'_a \subset Z \times_S S'$ を既約成分、その一般点を $y'_a \in Z'_a$ とし、
$Z \times_S S'$ における $Z'_a$ の重複度を $n_a$ とする。このとき
$$
\beta' = \sum [Z \times_S S']_{e + c} = \sum n_a[Z'_a]
$$
$\alpha'$ は $Y' \to Y$ による $\alpha$ の基底変換なので
$\alpha'_{y'_a} = \sum m_j [V_{j, \kappa(y'_a)}]_r$ である。
$V'_{jab} \subset V_{j, \kappa(y'_a)}$ を既約成分とし、
$V_{j, \kappa(y'_a)}$ における $V'_{jab}$ の重複度を $m_{jab}$ とする。このとき
$$
\alpha'_{y'_a} = \sum m_j [V_{j, \kappa(y'_a)}]_r =
\sum m_j m_{jab} [V'_{jab}]
$$
したがって
$$
\alpha' \cap \beta' = \sum n_a m_j m_{jab} [\overline{V}'_{jab}]
$$
である。ここで $\overline{V}'_{jab}$ は $X'$ における $V'_{jab}$ の閉包である。
したがって所望の等式を証明するには、次を示せば十分である：
\begin{enumerate}
\item $\overline{V}_j \times_S S'$ の既約成分はスキーム
$\overline{V}'_{jab}$ である；
\item $\overline{V}_j \times_S S'$ における $\overline{V}'_{jab}$ の重複度は
$n_a m_{jab}$ に等しい。
\end{enumerate}
$V_j \to \overline{V}_j$ は整スキームの双有理射であることに注意せよ。射
$V_j \times_S S' \to V_j$ および
$\overline{V}_j \times_S S' \to \overline{V}_j$ は平坦なので、既約成分の一般点を
$V_j$ および $\overline{V}_j$ の (一意な) 一般点へ写す。したがって
$V_j \times_S S' \to \overline{V}_j \times_S S'$ は双有理射であり、既約成分の間の
全単射を誘導してその重複度を同一視する。ゆえに、$V_j \times_S S'$ の既約成分が
$V'_{jab}$ であり、$V_j \times_S S'$ における $V'_{jab}$ の重複度が
$n_a m_{jab}$ に等しいことを示せば十分である。しかしこれは、図式
$$
\xymatrix{
Z_r(V_j) \ar[r] & Z_{r + c}(V_j \times_S S') \\
Z_0(\Spec(\kappa(y))) \ar[r] \ar[u] &
Z_c(\Spec(\kappa(y)) \times_S S') \ar[u]
}
$$
が可換であるということにほかならない。ここで横の射は
$\Spec(\kappa(y)) \times_S S' \to \Spec(\kappa(y))$ による基底変換、
縦の射は平坦引き戻しである。これは Chow Homology, Lemma
\ref{chow-lemma-pullback-base-change-pullback} で示されている。

\medskip\noindent
補題における平坦引き戻しに関する主張は、次のことを意味する。
$(S, \delta)$, $X \to Y$, $\alpha$, $\beta$ を、上の
$\alpha \cap \beta$ の構成におけるものとする。$Y' \to Y$ を
局所有限型、相対次元 $c$ の平坦射とする。このとき $\alpha'$ を
$Y' \to Y$ による $\alpha$ の基底変換、$\beta'$ を $\beta$ の
平坦引き戻しとすることができる。平坦引き戻しとの両立性とは、
$\alpha' \cap \beta'$ が $X \times_Y Y' \to Y$ による
$\alpha \cap \beta$ の平坦引き戻しであることを意味する。実際これは、
$S = Y$, $S' = Y'$ とおけば、上の議論の特別な場合である。
\end{proof}

\begin{lemma}
\label{relative-cycles-lemma-action-coherent}
$(S, \delta)$ および $f : X \to Y$ を上のようなものとする。
$\mathcal{F}$ を、すべての $y \in Y$ に対して
$\dim(\text{Supp}(\mathcal{F}_y)) \leq r$ を満たす連接
$\mathcal{O}_X$-加群とする。$\mathcal{G}$ を
$\dim_\delta(\text{Supp}(\mathcal{G})) \leq e$ を満たす連接
$\mathcal{O}_Y$-加群とする。$\alpha = [\mathcal{F}/X/Y]_r$
(Example \ref{relative-cycles-example-family-associated-module}) および
$\beta = [\mathcal{G}]_e$ (Chow Homology, Definition
\ref{chow-definition-cycle-associated-to-coherent-sheaf}) とおく。
$\mathcal{F}$ が $Y$ 上平坦ならば、
$\alpha \cap \beta =
[\mathcal{F} \otimes_{\mathcal{O}_X} f^*\mathcal{G}]_{r + e}$ である。
\end{lemma}

\begin{proof}
次の等式に注意する：
$$
\text{Supp}(\mathcal{F} \otimes_{\mathcal{O}_X} f^*\mathcal{G}) =
\text{Supp}(\mathcal{F}) \cap f^{-1}\text{Supp}(\mathcal{G}) =
\bigcup\nolimits_{y \in \text{Supp}(\mathcal{G})} \text{Supp}(\mathcal{F}_y)
$$
したがってこれは $\delta$-次元 $\leq r + e$ の閉部分集合である。ゆえに
$[\mathcal{F} \otimes_{\mathcal{O}_X} f^*\mathcal{G}]_{r + e}$
は定義される。

\medskip\noindent
以下では、$\alpha \cap \beta$ の構成で導入した記法
$\beta = \sum n_i[Z_i]$, $y_i \in Z_i$,
$\alpha_{y_i} = \sum m_{ij} [V_{ij}]$、および
$\overline{V}_{ij}$ を用いる。$\beta = [\mathcal{G}]_e$ なので、
$Z_i$ は $\text{Supp}(\mathcal{G})$ のうち $\delta$-次元が $e$ の
既約成分である。同様に、$V_{ij}$ は
$\text{Supp}(\mathcal{F}_{y_i})$ のうち次元が $r$ の既約成分である。
このことと第一段落の等式から、$\overline{V}_{ij}$ は
$\text{Supp}(\mathcal{F} \otimes_{\mathcal{O}_X} f^*\mathcal{G})$
のうち $\delta$-次元が $r + e$ の既約成分である。したがって補題を
証明するには、次を示せば十分である：
$$
\text{length}_{\mathcal{O}_{X, \xi_{ij}}}(
(\mathcal{F} \otimes_{\mathcal{O}_X} f^*\mathcal{G})_{\xi_{ij}})
=
\text{length}_{\mathcal{O}_{X_{y_i}, \xi_{ij}}}((\mathcal{F}_{y_i})_{\xi_{ij}})
\cdot
\text{length}_{\mathcal{O}_{Y, y_i}}(\mathcal{G}_{y_i})
$$
証明の第一段落により、左辺は $B = \mathcal{O}_{X, \xi_{ij}}$-加群
$$
\mathcal{G}_{y_i}
\otimes_{\mathcal{O}_{Y, y_i}}
\mathcal{F}_{\xi_{ij}} =
M \otimes_A N
$$
の長さに等しい。ここで $M = \mathcal{G}_{y_i}$ は有限長の
$A = \mathcal{O}_{Y, y_i}$-加群であり、$N = \mathcal{F}_{\xi_{ij}}$ は
$N/\mathfrak m_AN$ が有限長となる有限 $B$-加群である。
$\mathcal{F}$ は $Y$ 上平坦なので、$N$ は $A$-平坦である。式の右辺は
$$
\text{length}_B(N/\mathfrak m_A N) \cdot \text{length}_A(M)
$$
に等しい。したがって式の左辺と右辺はいずれも $M$ に関して加法的である
($N$ の $A$ 上の平坦性を用いる)。ゆえに $M = \kappa_A$ が剰余体である場合に
式を証明すれば十分であり、この場合は直ちに従う。
\end{proof}

\begin{lemma}
\label{relative-cycles-lemma-action-closed}
$(S, \delta)$ および $f : X \to Y$ を上のようなものとする。
$Z \subset X$ を $Y$ 上相対次元 $\leq r$ の閉部分スキームとし、
$\alpha = [Z/X/Y]_r$ (Example \ref{relative-cycles-example-family-associated-closed}) とおく。
$W \subset Y$ を $\delta$-次元 $\leq e$ の閉部分スキームとし、
$\beta = [W]_e$ (Chow Homology, Definition
\ref{chow-definition-cycle-associated-to-closed-subscheme}) とおく。
$Z$ が $Y$ 上平坦ならば、
$\alpha \cap \beta = [Z \times_Y W]_{r + e}$ である。
\end{lemma}

\begin{proof}
$\mathcal{F} = \mathcal{O}_Z$ および $\mathcal{F} = \mathcal{O}_W$ とおけば、
Lemma \ref{relative-cycles-lemma-action-coherent} の特別な場合である。
\end{proof}

\begin{lemma}
\label{relative-cycles-lemma-flat-pullback-as-action}
$(S, \delta)$ および $f : X \to Y$ を上のようなものとする。
$f$ が相対次元 $e$ の平坦射であると仮定する。$\beta \in Z_r(Y)$ に対し、
$Z_{e + r}(X)$ において $f^*\beta = [X/X/Y]_e \cap \beta$ が成り立つ。
\end{lemma}

\begin{proof}
等式は $X$ 上局所的に証明すれば十分である。したがって $Y$ はアフィンであると
仮定してよい。この場合、$\beta$ は素サイクルの有限整係数線形結合である。
等式の両辺は $\beta$ に関して加法的なので、$W$ が $Y$ の
$\delta$-次元 $r$ の整閉部分スキームで、$\beta = [W]$ であると仮定してよい。
$f^{-1}(W) = W \times_Y X$ を $X$ における $W$ のスキーム論的逆像とする。
Lemma \ref{relative-cycles-lemma-action-closed} により
$[X/X/Y]_e \cap \beta = [W \times_Y X]_{r + e}$ であり、
Chow Homology, Definition \ref{chow-definition-flat-pullback} により
$f^*\beta = [f^{-1}(W)]_{r + e}$ である。よって所望の等式を得る。
\end{proof}

\begin{lemma}
\label{relative-cycles-lemma-action-push-pull}
$(S, \delta)$ を上のようなものとする。
$$
\xymatrix{
X' \ar[r]_f \ar[d] & X \ar[d] \\
Y' \ar[r]^g & Y
}
$$
を $S$ 上局所有限型なスキームの Cartesian 図式とし、$g$ は固有であるとする。
$r, e \geq 0$ とする。$\alpha$ を $X/Y$ のファイバー上の
$r$-サイクルの族とし、$\beta' \in Z_e(Y')$ とする。このとき
$f_*(g^*\alpha \cap \beta') = \alpha \cap g_*\beta'$ が成り立つ。
\end{lemma}

\begin{proof}
$X$ 上のサイクルの等式を証明するので、Lemma \ref{relative-cycles-lemma-action-opens} により
$Y$ 上局所的に調べてよい。したがって $Y$ はアフィンであると仮定してよい。
すると $Y'$ は準コンパクトである。特に $\beta'$ は素サイクルの有限線形結合である。
$- \cap -$ は第二変数について線形なので (Lemma \ref{relative-cycles-lemma-action-bilinear})、
$Z' \subset Y'$ が $\delta$-次元 $e$ の整閉部分スキームで
$\beta' = [Z']$ である場合に等式を証明すれば十分である。$Z = g(Z')$ とおく。
これは $Y$ の $\delta$-次元 $\leq e$ の整閉部分スキームである。
簡単のため $Z$ の $\delta$-次元が $e$ に等しい場合を仮定し、他の場合
（こちらの方が容易である）は読者に委ねる。$y \in Z$ および $y' \in Z'$ を
一般点とする。$V_j \subset X_y$ を次元 $r$ の整閉部分スキームとして、
$\alpha_y = \sum m_j[V_j]$ と書く。

\medskip\noindent
まず $g$ が閉埋め込みであると仮定する。このとき $g_*\beta' = [Z]$ かつ
$(g^*\alpha)_{y'} = \sum n_j[V_j]$ である。後者は、$V_j$ が $X_y$ の
閉部分スキーム $X'_{y'}$ に含まれるので意味をもつ。この場合、等式は明らかである。
いずれの場合にも $\sum m_j[\overline{V}_j]$ を得る。ここで
$\overline{V}_j$ は閉部分スキーム $X' \subset X$ における $V_j$ の閉包である。

\medskip\noindent
上のように $\beta' = [Z']$ である一般の場合に戻る。
$W = Z \times_Y X$ および $W' = Z' \times_{Y'} X'$ とおく。Cartesian 正方形
$$
\xymatrix{
W \ar[r] \ar[d] & X \ar[d] \\
Z \ar[r] & Y
}
\quad
\xymatrix{
W' \ar[r] \ar[d] & X' \ar[d] \\
Z' \ar[r] & Y'
}
\quad
\xymatrix{
W' \ar[r] \ar[d] & W \ar[d] \\
Z' \ar[r] & Z
}
$$
を考える。前段落により最初の二つの正方形については結果が分かっているので、
形式的な議論により、最後の正方形と元 $\beta' = [Z'] \in Z_e(Z')$ に対して
結果を証明すれば十分である。これで次の段落で扱う場合に帰着される。

\medskip\noindent
$Y' \to Y$ が $\delta$-次元 $e$ の整スキームの生成的有限射であり、
$\beta' = [Y']$ であると仮定する。この場合、
$f_*(g^*\alpha \cap \beta')$ と $\alpha \cap g_*\beta'$ はともに、
$Y$ を支配する素サイクルの和として書けるサイクルである。したがって等式を
確かめるため、$Y$ を空でない開部分スキームで置き換えてよい。このように
置き換えた後、$g$ は次数 $d \geq 1$ の有限平坦射であると仮定してよい。
もちろん、これは $g_*\beta' = g_*[Y'] = d[Y]$ を意味する。また
$\beta' = [Y'] = g^*[Y]$ である。したがって
$$
f_*(g^*\alpha \cap \beta') =
f_*(g^*\alpha \cap g^*[Y]) =
f_*f^*(\alpha \cap [Y]) =
d (\alpha \cap [Y]) =
\alpha \cap g_*\beta'
$$
となり、所望の等式を得る。第二の等式は Lemma \ref{relative-cycles-lemma-action-base-change}、
第三の等式は Chow Homology, Lemma \ref{chow-lemma-finite-flat} による。
\end{proof}







\section{Chow 群への作用}
\label{relative-cycles-section-action-chow}

\noindent
$\alpha$ が相対 $r$-サイクルであるとき、Section \ref{relative-cycles-section-action} の作用素
$\alpha \cap -$ は有理同値を経由し、双変類を定める。

\begin{lemma}
\label{relative-cycles-lemma-closed-in-X-gysin}
$(S, \delta)$ を Section \ref{relative-cycles-section-action} のようなものとする。
$f : X' \to X$ を $S$ 上局所有限型なスキームの固有射とする。
$(\mathcal{L}, s, i : D \to X)$ を Chow Homology, Definition
\ref{chow-definition-gysin-homomorphism} のようなものとする。
Chow Homology, Remark \ref{chow-remark-pullback-pairs} のように図式
$$
\xymatrix{
D' \ar[d]_g \ar[r]_{i'} & X' \ar[d]^f \\
D \ar[r]^i & X
}
$$
を作る。$\mathcal{L}|_D \cong \mathcal{O}_D$ ならば、任意の
$\alpha' \in Z_{k + 1}(X')$ に対して $Z_k(D)$ において
$i^*f_*\alpha' = g_*(i')^*\alpha'$ が成り立つ。
\end{lemma}

\begin{proof}
すべての作用素はサイクルのレベルで定義されているので、この主張は意味をもつ。
Gysin 写像については Chow Homology, Remark
\ref{chow-remark-gysin-on-cycles} を参照せよ。ある整閉部分スキーム
$W' \subset X'$ に対して $\alpha = [W']$ であると仮定し、
$W = f(W') \subset X$ とおく。$W' \not \subset D'$ の場合、
$W \not \subset D$ であり、
$$
[W' \cap D']_k = \text{div}_{\mathcal{L}'|_{W'}}({s'|_{W'}})
\quad\text{and}\quad
[W \cap D]_k = \text{div}_{\mathcal{L}|_W}(s|_W)
$$
である。したがって Chow Homology, Lemma
\ref{chow-lemma-equal-c1-as-cycles} により、第一のサイクルの $f_*$ は
第二のサイクルに等しい。ゆえにサイクルとして等式が成り立つ。
$W' \subset D'$ の場合は $W \subset D$ であり、構成により両辺は零である。
\end{proof}

\begin{lemma}
\label{relative-cycles-lemma-action-gysin}
$(S, \delta)$ を Section \ref{relative-cycles-section-action} のようなものとする。
$X \to Y$ を $S$ 上局所有限型なスキームの射とする。$r \geq 0$ とし、
$\alpha \in z(X/Y, r)$ を $X/Y$ 上の相対 $r$-サイクルとする。
$(\mathcal{L}, s, i : D \to Y)$ を Chow Homology, Definition
\ref{chow-definition-gysin-homomorphism} のようなものとする。Cartesian 図式
$$
\xymatrix{
E \ar[d] \ar[r]_j & X \ar[d] \\
D \ar[r]^i & Y
}
$$
を作る。Chow Homology, Remark \ref{chow-remark-pullback-pairs} を参照せよ。
$\mathcal{L}|_D \cong \mathcal{O}_D$ ならば、$e \in \mathbf{Z}$ に対して図式
$$
\xymatrix{
Z_e(D) \ar[rr]_{i^*\alpha \cap -} & &
Z_{e + r}(E) \\
Z_{e + 1}(Y) \ar[u]^{i^*} \ar[rr]^{\alpha \cap -} & &
Z_{r + e + 1}(X) \ar[u]_{j^*}
}
$$
は可換である。ここで縦の射 $i^*$ および $j^*$ は、Chow Homology,
Remark \ref{chow-remark-gysin-on-cycles} のようなサイクル上の Gysin 写像である。
\end{lemma}

\begin{proof}
まず予備的な注意を述べる。$g : Y' \to Y$ を包絡射と仮定する
(Chow Homology, Definition \ref{chow-definition-envelope})。
$D, i, E, j, X, \alpha$ の $g$ による基底変換をそれぞれ
$D', i', E', j', X', \alpha'$ と書き、射影を $f : X' \to X$ と書く。
$D', i', E', j', X', Y', \alpha'$ に対して補題が成り立つと仮定する。
このとき $\beta' \in Z_{e + 1}(Y')$ ならば、
\begin{align*}
i^*\alpha \cap i^*g_*\beta'
& =
i^*\alpha \cap f_*(i')^*\beta' \\
& =
f_*(f^*i^*\alpha \cap (i')^*\beta') \\
& =
f_*((i')^*\alpha' \cap (i')^*\beta') \\
& =
f_*((j')^*(\alpha' \cap \beta')) \\
& =
j^*(f_*(f^*\alpha \cap \beta')) \\
& =
j^*(\alpha \cap g_*\beta')
\end{align*}
ここで第一の等式は Lemma \ref{relative-cycles-lemma-closed-in-X-gysin}、
第二の等式は Lemma \ref{relative-cycles-lemma-action-push-pull}、第三の等式は
$\alpha'$ の定義、第四の等式は $D', i', E', j', X', \alpha'$ に対して
本補題が成り立つという仮定、第五の等式は Lemma
\ref{relative-cycles-lemma-closed-in-X-gysin}、第六の等式は Lemma
\ref{relative-cycles-lemma-action-push-pull} による。したがって写像
$g_* : Z_{e + 1}(Y') \to Z_e(Y)$ の像に対して本補題が成り立つ。
しかし $g$ は完全分解的なので、この写像は全射である。ゆえに
$D, i, E, j, X, Y, \alpha$ に対して補題が成り立つ。

\medskip\noindent
$\beta \in Z_{e + 1}(Y)$ とする。$E$ 上のサイクルとして
$(D \to Y)^*\alpha \cap i^*\beta = j^*(\alpha \cap \beta)$
を示さなければならない。この問題は $E$ 上局所的なので、$X$ および $Y$ を
開部分スキームで置き換えてよい。（ここでは、作用素 $i^*$, $j^*$,
$\alpha \cap - $ および $(D \to Y)^*\alpha \cap -$ の形成が局所化と
可換であることを用いる。Gysin 写像については明らかであり、他については
Lemma \ref{relative-cycles-lemma-action-opens} から従う。）したがって $X$ と $Y$ は
アフィンであると仮定してよく、次の段落で扱う場合に帰着される。

\medskip\noindent
$X$ および $Y$ が準コンパクトであると仮定する。証明の第一段落と
Lemma \ref{relative-cycles-lemma-get-cycles} により、さらに $\alpha$ は
(\ref{relative-cycles-equation-cycle-classes}) の像に属すると仮定してよい。関係する作用素の
線形性により、$Y$ 上平坦かつ相対次元 $\leq r$ のある閉部分スキーム
$Z \subset X$ に対して $\alpha = [Z/X/Y]_r$ であると仮定してよい。
また $Y$ は準コンパクトなので、サイクル $\beta$ は素サイクルの有限線形結合である。
関係する作用素は線形だから、$W \subset Y$ が $\delta$-次元 $e + 1$ の
整閉部分スキームで $\beta = [W]$ である場合に等式を証明すれば十分である。

\medskip\noindent
$W \subset D$ ならば、一方で $i^*[W] = 0$ であり、他方で
$\alpha \cap [W]$ は $E$ 上に台をもつので
$j^*(\alpha \cap [W]) = 0$ でもある。したがってこの場合、等式は成り立つ。

\medskip\noindent
$W \not \subset D$ とする。このとき $i^*[W] = [D \cap W]_e$ である。
$\alpha = [Z/X/Y]_r$ の $i$ による引き戻し $i^*\alpha$ は
$[(E \cap Z)/E/D]_r$ であり、
$(E \cap Z) = E \times_Y Z = D \times_Y Z$
は $D$ 上平坦であることに注意せよ。したがって Lemma
\ref{relative-cycles-lemma-action-closed} を二度用いると、
$$
i^*\alpha \cap i^*[W] =
[(E \cap Z) \times_D (D \cap W)]_{r + e} =
[E \cap (Z \times_Y W)]_{r + e} =
j^*(\alpha \cap [W])
$$
となり、所望の等式を得る。
\end{proof}

\begin{proposition}
\label{relative-cycles-proposition-get-bivariant-class}
$(S, \delta)$ を Section \ref{relative-cycles-section-action} のようなものとする。
$X \to Y$ を $S$ 上局所有限型なスキームの射とする。$r \geq 0$ とし、
$\alpha \in z(X/Y, r)$ を $X/Y$ 上の相対 $r$-サイクルとする。
各局所有限型射 $g : Y' \to Y$ と各 $e \in \mathbf{Z}$ に作用素
$$
g^*\alpha \cap - : Z_e(Y') \to Z_{r + e}(X')
$$
を対応させる規則（ここで $X' = Y' \times_Y X$ とする）は有理同値を経由し、
双変類 $c(\alpha) \in A^{-r}(X \to Y)$ を定める。
\end{proposition}

\begin{proof}
Lemma \ref{relative-cycles-lemma-action-gysin} と Chow Homology, Lemma
\ref{chow-lemma-factors-through-rational-equivalence} により、この作用素は
有理同値を経由する。Chow Homology, Lemma
\ref{chow-lemma-bivariant-weaker} および Lemmas
\ref{relative-cycles-lemma-action-push-pull}, \ref{relative-cycles-lemma-action-base-change},
\ref{relative-cycles-lemma-action-gysin} により、得られる Chow 群上の作用素は双変類である。
\end{proof}

\begin{remark}
\label{relative-cycles-remark-characterize-relative-cycles}
$(S, \delta)$ を Section \ref{relative-cycles-section-action} のようなものとする。
$X \to Y$ を $S$ 上局所有限型なスキームの射とし、$r \geq 0$ とする。
$c$ を、各局所有限型射 $g : Y' \to Y$ と各 $e \in \mathbf{Z}$ に作用素
$$
c \cap - : Z_e(Y') \to Z_{r + e}(X')
$$
を対応させる規則で、固有前送り、平坦引き戻し、および Lemma
\ref{relative-cycles-lemma-action-gysin} のような Gysin 写像と両立するものとする。このとき、
上の各 $g$ に対して $c \cap = g^*\alpha \cap -$ となるような $X/Y$ 上の
相対 $r$-サイクル $\alpha$ が存在すると主張する。この結果が必要になれば、
ここで注意深く定式化して証明することにする。
\end{remark}









\section{ファイバー上のサイクルの族の合成}
\label{relative-cycles-section-compose-families}

\noindent
$X \to Y \to S$ をともに局所有限型なスキームの射とする。
$r, e \geq 0$ とする。$\alpha$ を $X/Y$ のファイバー上の $r$-サイクルの族、
$\beta$ を $Y/S$ のファイバー上の $e$-サイクルの族とする。このとき
$$
(\alpha \circ \beta)_s = (Y_s \to Y)^*\alpha \cap \beta_s
$$
とおくことで、$X/S$ のファイバー上の $(r + e)$-サイクルの族
$\alpha \circ \beta$ を得る。より正確には、$(Y_s \to Y)^*\alpha$ は
$Y_s \to Y$ による $\alpha$ の基底変換、すなわち $X_s/Y_s$ のファイバー上の
$r$-サイクルの族を表す。また作用素 $- \cap -$ は Section
\ref{relative-cycles-section-action} で定義され、調べられたものである\footnote{念のため、
ここでは $s = \Spec(\kappa(s))$ を $\delta(s) = 0$ をもつ基底スキームとして用いる。}。

\begin{lemma}
\label{relative-cycles-lemma-construction-bilinear}
上の構成は双線形である。すなわち
$(\alpha_1 + \alpha_2) \circ \beta \alpha_1 \circ \beta +
\alpha_1 \circ \beta$ および $\alpha \circ (\beta_1 + \beta_2) =
\alpha \circ \beta_1 + \alpha \circ \beta_2$ が成り立つ。
\end{lemma}

\begin{proof}
省略する。ヒント：Lemma \ref{relative-cycles-lemma-action-bilinear} により、ファイバー上で
この構成は双線形である。
\end{proof}

\begin{lemma}
\label{relative-cycles-lemma-construction-opens}
$U \subset X$ および $V \subset Y$ が開で $f(U) \subset V$ ならば、
$(\alpha \circ \beta)|_U$ は $\alpha|_U \circ \beta|_V$ に等しい。
\end{lemma}

\begin{proof}
省略する。ヒント：ファイバー上で Lemma \ref{relative-cycles-lemma-action-opens} を用いよ。
\end{proof}

\begin{lemma}
\label{relative-cycles-lemma-construction-base-change}
$\alpha \circ \beta$ の形成は基底変換と両立する。
\end{lemma}

\begin{proof}
$g : S' \to S$ をスキームの射とする。$X \to Y$ の $g$ による基底変換を
$X' \to Y'$ と書く。$Y' \to Y$ に関する $\alpha$ の基底変換を $\alpha'$、
$S' \to S$ に関する $\beta$ の基底変換を $\beta'$ と書く。主張は、
$\alpha' \circ \beta'$ が $g : S' \to S$ による $\alpha \circ \beta$ の
基底変換であることを意味する。

\medskip\noindent
$s' \in S'$ を像 $s \in S$ をもつ点とする。このとき
$$
(\alpha' \circ \beta')_{s'} = (Y'_{s'} \to Y')^*\alpha' \cap \beta'_{s'}
$$
である。また
$$
(Y'_{s'} \to Y')^*\alpha' =
(Y'_{s'} \to Y')^*(Y' \to Y)^*\alpha =
(Y'_{s'} \to Y_s)^*(Y_s \to Y)^*\alpha
$$
であり、$\beta'_{s'}$ は
$s' = \Spec(\kappa(s')) \to \Spec(\kappa(s)) = s$ による $\beta_s$ の
基底変換である。したがって $(Y_s \to Y)^*\alpha$, $\beta_s$,
$X_s \to Y_s \to s$ および $s' \to s$ による基底変換に
Lemma \ref{relative-cycles-lemma-action-base-change} を適用すれば結果が従う。
\end{proof}

\begin{lemma}
\label{relative-cycles-lemma-construction-coherent}
$f : X \to Y$ および $Y \to S$ をともに局所有限型なスキームの射とする。
$r, e \geq 0$ とする。$\mathcal{F}$ を、すべての $y \in Y$ に対して
$\dim(\text{Supp}(\mathcal{F}_y)) \leq r$ を満たす有限型準連接
$\mathcal{O}_X$-加群とする。$\mathcal{G}$ を、すべての $s \in S$ に対して
$\dim(\text{Supp}(\mathcal{G}_s)) \leq e$ を満たす有限型準連接
$\mathcal{O}_Y$-加群とする。$\alpha = [\mathcal{F}/X/Y]_r$ および
$\beta = [\mathcal{G}/Y/S]_e$ とし
(Example \ref{relative-cycles-example-family-associated-module})、$\mathcal{F}$ が $Y$ 上平坦ならば、
$\alpha \circ \beta =
[\mathcal{F} \otimes_{\mathcal{O}_X} f^*\mathcal{G}/X/S]_{r + e}$ である。
\end{lemma}

\begin{proof}
まず $\mathcal{F} \otimes_{\mathcal{O}_X} f^*\mathcal{G}$ は有限型準連接
$\mathcal{O}_X$-加群である。$s \in S$ とする。テンソル積の右完全性により
$$
(\mathcal{F} \otimes_{\mathcal{O}_X} f^*\mathcal{G})_s =
\mathcal{F}_s \otimes_{\mathcal{O}_{X_s}} f_s^*\mathcal{G}_s
$$
である。さらに $\mathcal{F}_s$ は平坦加群の基底変換なので $Y_s$ 上平坦である。
したがって Lemma \ref{relative-cycles-lemma-action-coherent} により等式
$(\alpha \circ \beta)_s =
[(\mathcal{F} \otimes_{\mathcal{O}_X} f^*\mathcal{G})_s]_{r + e}$
が従う。
\end{proof}

\begin{lemma}
\label{relative-cycles-lemma-construction-closed}
$f : X \to Y$ および $Y \to S$ をともに局所有限型なスキームの射とする。
$r, e \geq 0$ とする。$Z \subset X$ を $Y$ 上相対次元 $\leq r$ の
閉部分スキームとする。$W \subset Y$ を $S$ 上相対次元 $\leq e$ の
閉部分スキームとする。$\alpha = [Z/X/Y]_r$ および $\beta = [W/Y/S]_e$ とし
(Example \ref{relative-cycles-example-family-associated-closed})、$Z$ が $Y$ 上平坦ならば、
$\alpha \circ \beta = [Z \times_Y W/X/S]_{r + e}$ である。
\end{lemma}

\begin{proof}
$\mathcal{F} = \mathcal{O}_Z$ および $\mathcal{F} = \mathcal{O}_W$ とおけば、
Lemma \ref{relative-cycles-lemma-construction-coherent} の特別な場合である。
\end{proof}

\begin{lemma}
\label{relative-cycles-lemma-flat-pullback-as-composition}
$f : X \to Y$ および $Y \to S$ をともに局所有限型なスキームの射とする。
$f$ が相対次元 $e$ の平坦射であると仮定する。$\beta$ を $Y/S$ の
ファイバー上の $r$-サイクルの族とする。このとき
$f^*\beta = [X/X/Y]_e \circ \beta$
が $X/S$ のファイバー上の $(e + r)$-サイクルの族として成り立つ。
\end{lemma}

\begin{proof}
定義を展開し、$[X/X/Y]_e$ の形成が基底変換と両立すること
(Lemma \ref{relative-cycles-lemma-family-associated-closed}) を用いると、これは
Lemma \ref{relative-cycles-lemma-flat-pullback-as-action} から従う。
\end{proof}

\begin{lemma}
\label{relative-cycles-lemma-construction-push-pull}
$S$ をスキームとする。
$$
\xymatrix{
X' \ar[r]_f \ar[d] & X \ar[d] \\
Y' \ar[r]^g & Y
}
$$
を $S$ 上局所有限型なスキームの Cartesian 図式とし、$g$ は固有であるとする。
$r, e \geq 0$ とする。$\alpha$ を $X/Y$ のファイバー上の $r$-サイクルの族、
$\beta'$ を $Y'/S$ のファイバー上の $e$-サイクルの族とする。このとき
$f_*(g^*(\alpha) \circ \beta') = \alpha \circ g_*\beta'$ が成り立つ。
\end{lemma}

\begin{proof}
定義を展開すれば、Lemma \ref{relative-cycles-lemma-action-push-pull} から従う。
\end{proof}

\begin{lemma}
\label{relative-cycles-lemma-construction-composition}
$(S, \delta)$ を Chow Homology, Situation \ref{chow-situation-setup} の
ようなものとする。$X \to Y \to Z$ を $S$ 上局所有限型なスキームの射とする。
$r, s, e \geq 0$ とする。このとき
$$
(\alpha \circ \beta) \cap \gamma = \alpha \cap (\beta \cap \gamma)
\quad\text{in}\quad Z_{r + s + e}(X)
$$
が成り立つ。ここで $\alpha$ は $X/Y$ のファイバー上の $r$-サイクルの族、
$\beta$ は $Y/Z$ のファイバー上の $s$-サイクルの族、
$\gamma \in Z_e(Z)$ である。
\end{lemma}

\begin{proof}
$X$ 上のサイクルの等式を証明するので、Lemma \ref{relative-cycles-lemma-action-opens} により
$Z$ 上局所的に調べてよい。したがって $Z$ はアフィンであると仮定してよい。
特に $\gamma$ は素サイクルの有限線形結合である。$- \cap -$ は第二変数について
線形なので (Lemma \ref{relative-cycles-lemma-action-bilinear})、$W \subset Z$ が
$\delta$-次元 $e$ の整閉部分スキームで $\gamma = [W]$ である場合に
等式を証明すれば十分である。

\medskip\noindent
$z \in W$ を一般点とし、$Z_s(Y_z)$ において
$\beta_z = \sum m_j[V_j]$ と書く。このとき $\beta \cap \gamma$ は
$\sum m_j[\overline{V}_j]$ に等しい。ここで $\overline{V}_j \subset Y$ は、
$Y \to Z$ により $W$ 内へ写され、一般ファイバーが $V_j$ である整閉部分スキームである。
$y_j \in V_j$ を一般点とする。これを $\overline{V}_j$ の一般点ともみなす
（これは $W$ の $z$ へ写る）。$Z_r(X_{y_j})$ において
$\alpha_{y_j} = \sum n_{jk} [W_{jk}]$ と書く。このとき
$\alpha \cap (\beta \cap \gamma)$ は
$$
\sum m_j n_{jk} [\overline{W}_{jk}]
$$
に等しい。ここで $\overline{W}_{jk} \subset X$ は、$X \to Y$ により
$\overline{V}_j$ 内へ写され、一般ファイバーが $W_{jk}$ である整閉部分スキームである。

\medskip\noindent
一方、
$$
(\alpha \circ \beta)_z = (Y_z \to Y)^*\alpha \cap \beta_z =
(Y_z \to Y)^*\alpha \cap (\sum m_j [V_j])
$$
を考える。$- \cap -$ の構成により、これは $X_z$ 上のサイクル
$$
\sum m_j n_{jk} [(\overline{W}_{jk})_z]
$$
に等しい。したがって定義により
$$
(\alpha \circ \beta) \cap [W] =
\sum m_j n_{jk} [\widetilde{W}_{jk}]
$$
を得る。ここで $\widetilde{W}_{jk} \subset X$ は、$X \to Z$ により
$W$ 内へ写され、一般ファイバーが $(\overline{W}_{jk})_z$ である整閉部分スキームである。
明らかに $\widetilde{W}_{jk} = \overline{W}_{jk}$ でなければならない。
これで証明が完了する。
\end{proof}









\section{相対サイクルの合成}
\label{relative-cycles-section-compose}

\noindent
$S$ を局所 Noether スキームとし、$X \to Y$ を $S$ 上局所有限型な
スキームの射とする。Section \ref{relative-cycles-section-compose-families} の構成を用いて写像
$$
z(X/Y, r) \otimes_\mathbf{Z} z(Y/S, e) \longrightarrow z(X/S, r + e),\quad
\alpha \otimes \beta \longmapsto \alpha \circ \beta
$$
を定義する。この構成はすでに双線形であることが分かっている
(Lemma \ref{relative-cycles-lemma-construction-bilinear})。したがって次を示せば、
表示された射を得る。

\begin{lemma}
\label{relative-cycles-lemma-well-defined}
$\alpha$ および $\beta$ が相対サイクルならば、$\alpha \circ \beta$ も
相対サイクルである。
\end{lemma}

\begin{proof}
Lemma \ref{relative-cycles-lemma-construction-base-change} により、$\alpha \circ \beta$ の形成は
基底変換と両立する。したがって $S$ は、一般点 $\eta$ と閉点 $0$ をもつ
離散付値環のスペクトルであると仮定してよく、
$sp_{X/S}((\alpha \circ \beta)_\eta) = (\alpha \circ \beta)_0$.
を示さなければならない。サイクルの等式を証明するので、$Y$ および $X$ 上
局所的に調べてよい（構成が制限と可換することを見るために Lemmas
\ref{relative-cycles-lemma-construction-opens} および
\ref{relative-cycles-lemma-specialization-flat-pullback} を用いる）。したがって $X$ と $Y$ は
アフィンであると仮定してよい。Lemma \ref{relative-cycles-lemma-get-cycles} により、
$g^*\alpha$ が (\ref{relative-cycles-equation-cycle-classes}) の像に属するような
完全分解的固有射 $g : Y' \to Y$ を見いだせる。

\medskip\noindent
射 $g_\eta : Y'_\eta \to Y_\eta$ は完全分解的なので、
$\beta_\eta = \sum g_{\eta, *}\beta'_\eta$
を満たす $\beta'_\eta \in Z_e(Y'_\eta)$ を見いだせる。Chow Homology,
Lemma \ref{chow-lemma-envelope} を参照せよ。
$\beta'_0 = sp_{Y'/S}(\beta'_\eta)$ とおくと、
$\beta' = (\beta'_\eta, \beta'_0)$ は $Y'/S$ 上の相対 $e$-サイクルである。
このとき $g_*\beta'$ および $\beta$ は $Y/S$ 上の相対 $e$-サイクルであり
(Lemma \ref{relative-cycles-lemma-relative-cycle-functoriality})、$\eta$ における値が同じなので
等しい (Lemma \ref{relative-cycles-lemma-uniqueness-extension})。線形性
(Lemma \ref{relative-cycles-lemma-construction-bilinear}) により、
$\alpha \circ g_*\beta'$ が相対 $(r + e)$-サイクルであることを示せば十分である。

\medskip\noindent
$X' = X \times_Y Y'$ とおき、射影を $f : X' \to X$ と書く。
Lemma \ref{relative-cycles-lemma-construction-push-pull} により
$\alpha \circ g_*\beta' = f_*(g^*\alpha \circ \beta')$.
である。Lemma \ref{relative-cycles-lemma-relative-cycle-functoriality} により、
$g^*\alpha \circ \beta'$ が相対 $(r + e)$-サイクルであることを示せば十分である。
Lemma \ref{relative-cycles-lemma-get-cycles-dvr} と双線形性を用いると、次の段落で扱う場合に帰着される。

\medskip\noindent
$\alpha = [Z/X/Y]_r$ および $\beta = [W/Y/S]$ と仮定する。ここで
$Z \subset X$ は $Y$ 上平坦かつ相対次元 $\leq r$ の閉部分スキームであり、
$W \subset Y$ は $S$ 上平坦かつ相対次元 $\leq e$ の閉部分スキームである。
Lemma \ref{relative-cycles-lemma-construction-closed} により
$$
\alpha \circ \beta = [Z \times_X W/X/S]_{r + e}
$$
であり、$Z \times_X W \subset X$ は $S$ 上平坦かつ相対次元
$\leq r + e$ の閉部分スキームである。Lemma
\ref{relative-cycles-lemma-family-associated-closed-specialization} により、これは
相対 $(r + e)$-サイクルである。
\end{proof}

\begin{lemma}
\label{relative-cycles-lemma-composition-fundamental-cycles}
$f : X \to Y$ および $g : Y \to S$ をスキームの射とする。
$S$ は局所 Noether、$g$ は局所有限型かつ相対次元 $e \ge 0$ の平坦射、
$f$ は局所有限型かつ相対次元 $r \geq 0$ の平坦射であると仮定する。
このとき $z(X/S, r + e)$ において
$[X/X/Y]_r \circ [Y/Y/S]_e = [X/X/S]_{r + e}$ が成り立つ。
\end{lemma}

\begin{proof}
Lemma \ref{relative-cycles-lemma-construction-closed} の特別な場合である。
\end{proof}






\section{Suslin および Voevodsky との比較}
\label{relative-cycles-section-compare}

\noindent
サイクルに関する記法を Chow Homology, Section \ref{chow-section-cycles}
以下から採ったことを除き、\cite{SV} と同じ記法を用いるよう努めた。
比較は次の通りである：
\begin{enumerate}
\item \cite[Section 3.1]{SV} には ``相対サイクル''、``次元 $r$ の相対サイクル''、
および ``次元 $r$ の等次元相対サイクル'' という概念がある。本章には
これらに対応する概念はない。したがって群
$Cycl(X/S, r)$, $Cycl_{equi}(X/S, r)$,
$PropCycl(X/S, r)$、および $PropCycl_{equi}(X/S, r)$
には本章に対応物がない。
\item \cite[page 36]{SV} の下部で群
$z(X/S, r)$, $c(X/S, r)$, $z_{equi}(X/S, r)$, $c_{equi}(X/S, r)$
が定義されている。$S$ が分離 Noether で、$X \to S$ が分離かつ有限型であるとき、
これらは本章の概念と一致する。
\item \cite{SV} では、記号 $z(X/S, r)$ が $S$ 上有限型なスキームの圏上の前層
$S' \mapsto z(S' \times_S X/S', r)$ を表すために用いられることがある。
$c(X/S, r)$, $z_{equi}(X/S, r)$, $c_{equi}(X/S, r)$ についても同様である。
\item \cite{SV} で定義される基底変換、平坦引き戻し、および固有前送りは、
双方が適用できる場合に本章のものと一致する。
\item $\alpha \in z(X/S, r)$ に対して Section \ref{relative-cycles-section-action} で定義された作用素
$\alpha \cap - : Z_e(S) \to Z_{e + r}(X)$ は、双方が定義される場合、
\cite[Section 3.7]{SV} の作用素 $Cor(\alpha, -)$ と一致する。
\item $X \to Y \to S$ に対して、Section \ref{relative-cycles-section-compose} で定義された合成則
$z(X/Y, r) \otimes_\mathbf{Z} z(Y/S, e) \longrightarrow z(X/S, r + e)$
は \cite[Corollary 3.7.5]{SV} の作用素 $Cor_{X/Y}(-, -)$ と一致する。
\end{enumerate}











\section{非 Noether の場合の相対サイクル}
\label{relative-cycles-section-non-noetherian}

\noindent
読者にはこの節を読み飛ばすことを強く勧める。

\medskip\noindent
$f : X \to S$ を有限表示なスキームの射とし、$r \geq 0$ とする。
$Z \to S$ が平坦、有限表示、かつ相対次元 $\leq r$ となる閉部分スキーム
$Z \subset X$ の集合を $Hilb(X/S, r)$ と書く。群準同型
\begin{equation}
\label{relative-cycles-equation-cycle-classes-general}
\begin{matrix}
\text{自由アーベル群} \\
\text{生成元 }Hilb(X/S, r)
\end{matrix}
\longrightarrow
\begin{matrix}
\text{ファイバー上の }r\text{-サイクル}\\
\text{基底 }X/S
\end{matrix}
\end{equation}
で、$\sum n_i[Z_i]$ を $\sum n_i[Z_i/X/S]_r$ へ写すものを考える。

\begin{lemma}
\label{relative-cycles-lemma-relative-r-cycle-general}
$S$ を準コンパクトかつ準分離なスキームとする。
$f : X \to S$ を有限表示な射とする。$r \geq 0$ とし、$\alpha$ を
$X/S$ のファイバー上の $r$-サイクルの族とする。次は同値である：
\begin{enumerate}
\item Cartesian 図式
$$
\xymatrix{
X \ar[r] \ar[d] & X_0 \ar[d] \\
S \ar[r] & S_0
}
$$
が存在し、$X_0 \to S_0$ は Noether スキームの有限型射で、
$\alpha_0 \in z(X_0/S_0, r)$ が存在して、$\alpha$ は $S \to S_0$ による
$\alpha_0$ の基底変換である；
\item $g^*\alpha$ が (\ref{relative-cycles-equation-cycle-classes-general}) の像に属するような、
有限表示な完全分解的固有射 $g : S' \to S$ が存在する。
\end{enumerate}
\end{lemma}

\begin{proof}
(1) のような図式と $\alpha_0 \in z(X_0/S_0, r)$ が与えられたとする。
Lemma \ref{relative-cycles-lemma-get-cycles} により、$g_0^*\alpha_0$ が
(\ref{relative-cycles-equation-cycle-classes-general}) の像に属するような完全分解的固有射
$g_0 : S'_0 \to S_0$ が存在する。実際、$S'_0$ は Noether なので、
$S'_0 \times_{S_0} X_0$ のすべての閉部分スキームは $S'_0$ 上有限表示である。
$S' = S \times_{S_0} S'_0$ とおき、$S' \to S'_0$ による基底変換を用いると
(2) が成り立つことが分かる。

\medskip\noindent
逆に (2) が成り立つと仮定する。$g^*\alpha$ が
(\ref{relative-cycles-equation-cycle-classes-general}) の像に属するような有限表示な
完全分解的固有射 $g : S' \to S$ を選ぶ。$X' = S' \times_S X$ とおく。
$Z_a \subset X'$ を $S'$ 上平坦、有限表示、かつ相対次元 $\leq r$ の
閉部分スキームとして、$g^*\alpha = \sum n_a [Z_a/X'/S']_r$ と書く。

\medskip\noindent
各 $S_i$ が $\mathbf{Z}$ 上有限型で、遷移射がアフィンとなる有向逆極限として
$S = \lim S_i$ と書く。Limits, Proposition \ref{limits-proposition-approximate}
を参照せよ。十分大きい $i$ を選べば、次が存在する：
\begin{enumerate}
\item $S$ への基底変換が $g : S' \to S$ となる完全分解的固有射
$g_i : S'_i \to S_i$；
\item $X'_i = S'_i \times_{S_i} X_i$ とおいたとき、$S'_i$ 上平坦かつ
相対次元 $\leq r$ で、$S'$ への基底変換が $Z_a$ となる閉部分スキーム
$Z_{ai} \subset X'_i$。
\end{enumerate}
このために Limits, Lemmas
\ref{limits-lemma-descend-finite-presentation},
\ref{limits-lemma-descend-closed-immersion-finite-presentation},
\ref{limits-lemma-descend-flat-finite-presentation},
\ref{limits-lemma-descend-surjective},
\ref{limits-lemma-eventually-proper}、および
\ref{limits-lemma-limit-dimension}
および More on Morphisms, Lemma \ref{more-morphisms-lemma-descend-cd} を用いる。
$\alpha'_i = \sum n_a [Z_{ai}/X'_i/S_i]_r \in z(X'_i/S'_i, r)$ を考える。
構成により、$\alpha'_i$ を $X'/S'$ のファイバー上の $r$-サイクルの族へ
基底変換したものは、基底変換 $g^*\alpha$ と一致する。

\medskip\noindent
$S''_i = S'_i \times_{S_i} S'_i$, $X''_i = S''_i \times_{S_i} X_i$ とおき、
$S'' = S' \times_S S'$, $X'' = S'' \times_S X$ とおく。射影を
$\text{pr}_1, \text{pr}_2 : S'' \to S'$ および
$\text{pr}_1, \text{pr}_2 : S''_i \to S'_i$ と書く。
$X''_i/S''_i$ 上の相対 $r$-サイクル $\text{pr}_1^*\alpha'_i$ および
$\text{pr}_1^*\alpha'_i$ は、
$\text{pr}_1^*g^*\alpha = \text{pr}_1^*g^*\alpha$.
であるため、$X''/S''$ のファイバー上の同じ $r$-サイクルの族へ基底変換される。
したがって射 $S'' \to S''_i$ の像は
$E =
\{s \in S''_i : (\text{pr}_1^*\alpha'_i)_s = (\text{pr}_1^*\alpha'_i)_s\}$.
に含まれる。Lemma \ref{relative-cycles-lemma-relative-cycles-equal} により、これは閉部分集合である。
$S'' = \lim_{i' \geq i} S''_{i'}$ なので、Limits, Lemma
\ref{limits-lemma-limit-contained-in-constructible} により、ある $i' \geq i$ に対して
射 $S''_{i'} \to S''_i$ の像は $E$ に含まれる。したがって $i$ を $i'$ で
置き換えた後、
$\text{pr}_1^*\alpha'_i = \text{pr}_1^*\alpha'_i$.
と仮定してよい。Lemma \ref{relative-cycles-lemma-descend-family} により、
$g_i^*\alpha_i = \alpha'_i$ を満たす $X_i/S_i$ のファイバー上の
$r$-サイクルの族 $\alpha_i$ を一意的に得る（ここでは $S'_i \to S_i$ が
完全分解的であることを用いる）。Lemma \ref{relative-cycles-lemma-relative-cycles-h-descent} により
$\alpha_i \in z(X_i/S_i, r)$ である。Lemma \ref{relative-cycles-lemma-descend-family} における
一意性から $\alpha_i$ の基底変換は $\alpha$ である。よって (1) が成り立つ。
\end{proof}

\noindent
{\bf 議論。}
$f : X \to S$, $r$, $\alpha$ を Lemma
\ref{relative-cycles-lemma-relative-r-cycle-general} のようなものとする。このとき
Lemma \ref{relative-cycles-lemma-relative-r-cycle-general} の同値な条件 (1), (2) が成り立つならば、
$\alpha$ は {\it $X/S$ 上の相対 $r$-サイクル}であると言うことができる。
この定義には多くの良い性質がある。例えば $S$ が Noether である場合には
以前の定義と矛盾せず、Section \ref{relative-cycles-section-families-specialization} の結果の
大部分はこの設定へ一般化される。

\medskip\noindent
さらに次のように一般化できる。$S$ を任意のスキームとし、
$f : X \to S$ を局所有限表示な射と仮定する。$r \geq 0$ とし、$\alpha$ を
$X/S$ のファイバー上の $r$-サイクルの族 $\alpha$ とする。このとき、
$f(U) \subset V$ を満たすアフィン開部分集合 $U \subset X$, $V \subset S$ に対して、
制限 $\alpha|_U$ が前段落で定義した $U/V$ 上の相対 $r$-サイクルであるならば、
$\alpha$ を {\it $X/S$ 上の相対 $r$-サイクル}という。この設定にも以前の結果の
多くが一般化される。

\medskip\noindent
これらの一般化が必要になれば、ここで注意深く定式化して証明することにする。
